% Computer Science A --- TypeScript (language fundamentals).

\runningheader{Computer Science A}{Web: TypeScript}{Page \thepage\ of \numpages}

\begingradingrange{csawebts}

\uplevel{\par\textbf{Part 1: What is TypeScript?}\par}

\question[4]
TypeScript is:
\begin{choices}
  \choice A browser plugin that blocks JavaScript execution
  \CorrectChoice An open-source language that adds optional static types on top of JavaScript
  \choice A replacement for the DOM API
  \choice A database query language for SQL servers only
\end{choices}
\begin{solution}
  Microsoft created and maintains TypeScript. It is widely used in Angular, many enterprise front-end codebases, and Node-based tooling. Source files usually end in \texttt{.ts} (or \texttt{.tsx} with JSX).
\end{solution}

\question[4]
Browsers run TypeScript source files \textbf{directly} without a build step in typical production apps.
\begin{choices}
  \choice True---modern browsers execute \texttt{.ts} natively
  \CorrectChoice False---TypeScript is compiled (transpiled) to JavaScript first, then the browser runs the \texttt{.js} output
  \choice True only when \texttt{strict} mode is off
  \choice False because TypeScript runs only on SQL servers
\end{choices}
\begin{solution}
  The TypeScript compiler (\texttt{tsc}) or a bundler (Webpack, Vite, esbuild) removes types and emits JavaScript. Types exist for developers and the checker---they are erased at compile time.
\end{solution}

\question[4]
Which features are added by TypeScript but \textbf{not} part of standard JavaScript syntax at runtime?
\begin{choices}
  \choice \texttt{function}, \texttt{return}, and \texttt{if}
  \CorrectChoice Interfaces, type annotations, enums (as TypeScript constructs), and many compile-time-only checks
  \choice \texttt{console.log} and \texttt{JSON.parse}
  \choice HTML template bindings
\end{choices}
\begin{solution}
  TypeScript is a \textbf{superset}: valid JavaScript is valid TypeScript (with few exceptions). Extra syntax helps catch mistakes before deployment---undefined property access, wrong argument types, unsafe \texttt{null} use.
\end{solution}

\question[4]
Why do large projects (including Angular) adopt TypeScript?
\begin{choices}
  \choice To eliminate HTML and CSS entirely
  \CorrectChoice For earlier error detection, clearer APIs via types, and better editor support (autocomplete, refactors)
  \choice Because browsers cannot run JavaScript without it
  \choice To avoid using \texttt{package.json}
\end{choices}
\begin{solution}
  Teams trade a compile step for safer refactors and documentation embedded in code. IDEs use the type checker to suggest members and flag broken calls while you type.
\end{solution}

\question[4]
A file named \texttt{app.component.ts} in an Angular project contains:
\begin{choices}
  \choice Only CSS rules
  \CorrectChoice TypeScript code (classes, decorators, types) that compiles to JavaScript for the browser
  \choice SQL migration scripts only
  \choice Raw HTML with no logic
\end{choices}
\begin{solution}
  Angular schematics generate \texttt{.ts} component and service files, \texttt{.html} templates, and \texttt{.css} styles. The build pipeline compiles TypeScript as part of \texttt{ng build}.
\end{solution}

\uplevel{\par\textbf{Part 2: JavaScript vs TypeScript}\par}

\question[4]
The relationship between JavaScript and TypeScript is best stated as:
\begin{choices}
  \choice TypeScript is unrelated to JavaScript and uses a different runtime
  \CorrectChoice TypeScript is a typed superset of JavaScript---TypeScript code includes JS plus optional type syntax
  \choice JavaScript is a subset of TypeScript that cannot run in browsers
  \choice They compile to Python
\end{choices}
\begin{solution}
  You can migrate gradually: rename \texttt{.js} to \texttt{.ts}, add types where helpful, and use \lstinline[style=tsinline]|// @ts-check| in plain JS files for lighter checking.
\end{solution}

\question[4]
When TypeScript reports a \textbf{type error}, that usually means:
\begin{choices}
  \choice The browser will show the same error to every user at runtime
  \CorrectChoice The program failed static analysis before run---fix the source or adjust types before shipping
  \choice \texttt{npm} is not installed
  \choice The DOM is missing a \texttt{<router-outlet>}
\end{choices}
\begin{solution}
  JavaScript often throws type-related bugs only when a branch runs (\texttt{undefined is not a function}). TypeScript catches many of those during \texttt{tsc} or IDE checks. Runtime logic errors still require tests and debugging.
\end{solution}

\question[4]
Given \lstinline[style=tsinline]|let id: string = 'a1';| followed later by \lstinline[style=tsinline]|id = 99;|:
\begin{choices}
  \choice TypeScript always allows this because numbers coerce to strings at compile time
  \CorrectChoice TypeScript reports a type error (\texttt{number} is not assignable to \texttt{string}) before the code runs
  \choice Only ESLint can detect this; \texttt{tsc} ignores it
  \choice The error appears only in CSS files
\end{choices}
\begin{solution}
  Plain JavaScript would allow the reassignment at runtime. TypeScript's checker rejects incompatible assignments during compilation. Similar checks catch misspelled property names on typed objects.
\end{solution}

\question[4]
After compilation, TypeScript \textbf{interfaces} in source code:
\begin{choices}
  \choice Become global objects in the browser
  \CorrectChoice Are removed---they do not exist in the emitted JavaScript
  \choice Replace all \texttt{function} keywords with \texttt{class}
  \choice Force Node.js to restart
\end{choices}
\begin{solution}
  Interfaces and most type aliases are compile-time only. Enums and some decorators may emit runtime code depending on settings; interfaces never do.
\end{solution}

\question[4]
Compared with JavaScript, TypeScript \textbf{variables} can:
\begin{choices}
  \choice Only hold numbers
  \CorrectChoice Be annotated with types (\lstinline[style=tsinline]|let n: number|) or inferred from initializers
  \choice Never use \lstinline[style=tsinline]|const| or \lstinline[style=tsinline]|let|
  \choice Not be passed to functions
\end{choices}
\begin{solution}
  JavaScript is dynamically typed at runtime. TypeScript adds optional static types so tools know intended shapes before execution.
\end{solution}

\question[4]
Which task still requires understanding \textbf{JavaScript} when using TypeScript?
\begin{choices}
  \choice None---TypeScript replaces JS knowledge entirely
  \CorrectChoice All of them: closures, promises, the event loop, DOM APIs, and npm/Node workflows still apply
  \choice Only CSS layout
  \choice Only SQL joins
\end{choices}
\begin{solution}
  TypeScript is JavaScript plus types. Learning TS without JS fundamentals leads to confusion at runtime when compiled code behaves like ordinary JS.
\end{solution}

\uplevel{\par\textbf{Part 3: TypeScript compiler and installation}\par}

\question[4]
The TypeScript compiler is invoked from the command line as:
\begin{choices}
  \choice \texttt{node compile.ts}
  \CorrectChoice \texttt{tsc} (often via \texttt{npx tsc} or an \texttt{npm run build} script)
  \choice \texttt{ng only} with no TypeScript package
  \choice \texttt{java -jar typescript.jar}
\end{choices}
\begin{solution}
  Install locally: \texttt{npm install typescript --save-dev}. Run \texttt{npx tsc} to use the project version. Global install (\texttt{npm install -g typescript}) is possible but teams prefer local versions pinned in \texttt{package.json}.
\end{solution}

\question[4]
The file \texttt{tsconfig.json} is used to:
\begin{choices}
  \choice Store Angular HTML templates
  \CorrectChoice Configure compiler options (target ECMAScript version, \texttt{strict}, included files, output directory)
  \choice Replace \texttt{package.json}
  \choice List Karma test matchers only
\end{choices}
\begin{solution}
  Common options: \texttt{compilerOptions.target}, \texttt{module}, \texttt{strict}, \texttt{outDir}, \texttt{rootDir}, \texttt{include}/\texttt{exclude}. \texttt{tsc} with no arguments reads \texttt{tsconfig.json} in the current folder.
\end{solution}

\question[4]
Running \texttt{tsc} on a \texttt{.ts} file typically produces:
\begin{choices}
  \choice A \texttt{.css} file only
  \CorrectChoice A \texttt{.js} file (and optionally \texttt{.d.ts} declaration files) with types stripped out
  \choice An executable binary for Windows only
  \choice A SQL database schema
\end{choices}
\begin{solution}
  Example: \texttt{greet.ts} $\rightarrow$ \texttt{greet.js}. Declaration files describe shapes for libraries consumed by other TS projects.
\end{solution}

\question[4]
\texttt{npm install typescript --save-dev} places TypeScript in:
\begin{choices}
  \choice \texttt{dependencies} (required for browser download at runtime)
  \CorrectChoice \texttt{devDependencies} (needed for build/test on developer machines, not shipped as app runtime code)
  \choice \texttt{bootstrap} array of \texttt{@NgModule}
  \choice \texttt{angular.json} \texttt{styles} only
\end{choices}
\begin{solution}
  Compilers, linters, and test runners are dev dependencies. The browser bundle contains compiled JavaScript, not the \texttt{tsc} binary.
\end{solution}

\question[4]
\texttt{npx tsc --init} creates:
\begin{choices}
  \choice A new Angular component
  \CorrectChoice A starter \texttt{tsconfig.json} with common defaults
  \choice A \texttt{package-lock.json} only
  \choice A production web server on port \texttt{4200}
\end{choices}
\begin{solution}
  Teams edit the generated config for \texttt{strict} checks, path aliases (\texttt{paths}), and which folders to compile (\texttt{src/**/*}).
\end{solution}

\question[4]
Watch mode recompiles when files change. The usual flag is:
\begin{choices}
  \choice \texttt{tsc --serve}
  \CorrectChoice \texttt{tsc --watch} (or \texttt{tsc -w})
  \choice \texttt{tsc --router}
  \choice \texttt{tsc --karma}
\end{choices}
\begin{solution}
  During development, \texttt{tsc -w} keeps output JS in sync. Angular CLI wraps similar tooling internally when you run \texttt{ng serve}.
\end{solution}

\newpage

\uplevel{\par\textbf{Part 4: \texttt{package.json}}\par}

\question[4]
In a Node/npm project, \texttt{package.json} primarily:
\begin{choices}
  \choice Stores compiled JavaScript bytecode only
  \CorrectChoice Describes the project metadata, dependencies, and script commands for npm
  \choice Replaces \texttt{tsconfig.json} for every compiler option
  \choice Is executed directly by Chrome without a build step
\end{choices}
\begin{solution}
  Fields include \texttt{name}, \texttt{version}, \texttt{scripts}, \texttt{dependencies}, and \texttt{devDependencies}. npm reads this file when you run \texttt{npm install} or \texttt{npm run <script>}.
\end{solution}

\question[4]
The \texttt{scripts} section in \texttt{package.json} lets you define:
\begin{choices}
  \choice CSS selectors for components
  \CorrectChoice Short commands such as \texttt{"build": "tsc"} run via \texttt{npm run build}
  \choice TypeScript interfaces
  \choice HTML template paths only
\end{choices}
\begin{solution}
\begin{lstlisting}[style=jsonexam]
{
  "scripts": {
    "build": "tsc",
    "start": "node dist/index.js"
  }
}
\end{lstlisting}
  Scripts standardize tasks for the whole team and for CI pipelines.
\end{solution}

\question[4]
Packages listed under \texttt{dependencies} are:
\begin{choices}
  \choice Only used on the developer laptop during compilation, never in production
  \CorrectChoice Required for the application to run in production (or consumed as a published library)
  \choice Always the same as \texttt{devDependencies}
  \choice Ignored by \texttt{npm install}
\end{choices}
\begin{solution}
  Example runtime deps: \texttt{@angular/core}, \texttt{rxjs}. Build-time tools (\texttt{typescript}, \texttt{karma}) belong in \texttt{devDependencies} for apps; libraries document what consumers need.
\end{solution}

\question[4]
After cloning a TypeScript project, the usual first step to install listed packages is:
\begin{choices}
  \choice \texttt{tsc --init} only
  \CorrectChoice \texttt{npm install} (reads \texttt{package.json} and creates/updates \texttt{node\_modules})
  \choice \texttt{ng serve} without installing anything
  \choice Delete \texttt{package-lock.json} before every build
\end{choices}
\begin{solution}
  \texttt{package-lock.json} (or \texttt{npm-shrinkwrap}) pins exact versions for reproducible installs. Commit the lockfile in team projects.
\end{solution}

\question[4]
The \texttt{main} field in \texttt{package.json} often indicates:
\begin{choices}
  \choice The HTML file Angular bootstraps
  \CorrectChoice The entry JavaScript file when another package imports your library
  \choice The CSS theme color
  \choice The Karma port number
\end{choices}
\begin{solution}
  For apps, \texttt{scripts.start} may run \texttt{node dist/index.js}. For libraries, \texttt{main} points to the compiled entry (for example \texttt{dist/index.js}) after \texttt{tsc} builds.
\end{solution}

\question[4]
Semantic versioning in \texttt{"version": "1.2.3"} refers to:
\begin{choices}
  \choice Hours, minutes, seconds of the build
  \CorrectChoice Major.minor.patch (breaking changes, features, fixes)
  \choice The number of TypeScript errors allowed
  \choice The count of files in \texttt{node\_modules}
\end{choices}
\begin{solution}
  Teams bump major version when breaking public APIs, minor for backward-compatible features, patch for bug fixes. npm ranges (\texttt{\^{}1.2.0}, \texttt{\~{}1.2.0}) control acceptable upgrades.
\end{solution}

\uplevel{\par\textbf{Part 5: Node.js as the JavaScript runtime}\par}

\question[4]
\textbf{Node.js} is:
\begin{choices}
  \choice A replacement for HTML in Angular templates
  \CorrectChoice A JavaScript runtime built on Chrome's V8 engine that runs JS outside the browser
  \choice The same program as the TypeScript compiler
  \choice A CSS framework from Microsoft
\end{choices}
\begin{solution}
  Node runs scripts on servers, developer machines, and CI agents---file I/O, HTTP servers, build tools. It is \textbf{not} required inside end users' browsers for a typical Angular SPA.
\end{solution}

\question[4]
When you run \texttt{npm install} or \texttt{npx tsc}, you are using:
\begin{choices}
  \choice The browser's rendering engine only
  \CorrectChoice Node.js to execute npm and the TypeScript compiler tooling on your machine
  \choice SQL Server
  \choice Angular's \lstinline[style=htmlexam]|<router-outlet>|
\end{choices}
\begin{solution}
  npm (Node Package Manager) ships with Node. Front-end developers install Node to get npm, run dev servers, and compile TypeScript even though users only need the compiled assets in the browser.
\end{solution}

\question[4]
TypeScript and Node.js are related because:
\begin{choices}
  \choice Node executes \texttt{.ts} files natively without compilation in all production setups
  \CorrectChoice Node runs the \texttt{tsc} tool and build scripts; TypeScript describes JS that may run in Node or in browsers after compile
  \choice They are the same product with different names
  \choice Node replaces \texttt{package.json}
\end{choices}
\begin{solution}
  You can write server apps in TypeScript (\texttt{ts-node}, or compile to JS first). Angular CLI uses Node for \texttt{ng serve}, \texttt{ng build}, and tests---then ships browser-ready JavaScript bundles.
\end{solution}

\question[4]
The folder \texttt{node\_modules} after \texttt{npm install} contains:
\begin{choices}
  \choice Only TypeScript interface files used by the checker
  \CorrectChoice Downloaded packages listed in \texttt{package.json} (and their dependencies)
  \choice Compiled Angular components ready for production
  \choice User-uploaded images from the browser
\end{choices}
\begin{solution}
  Do not edit files inside \texttt{node\_modules} by hand. Add dependencies with \texttt{npm install <pkg>}. Many projects git-ignore \texttt{node\_modules} because it is reproducible from the lockfile.
\end{solution}

\question[4]
\texttt{node app.js} differs from opening \texttt{app.html} in a browser because:
\begin{choices}
  \choice \texttt{node} renders HTML and CSS with the DOM API only
  \CorrectChoice \texttt{node} runs JavaScript in a server/desktop process without the browser document, \texttt{window}, or DOM (unless polyfilled)
  \choice \texttt{node} cannot read files
  \choice \texttt{node} only runs TypeScript, not JavaScript
\end{choices}
\begin{solution}
  Node provides APIs such as \texttt{fs}, \texttt{http}, and \texttt{path}. Front-end code expecting \texttt{document} must run in a browser or a DOM test environment (JSDOM).
\end{solution}

\question[4]
For an Angular student workstation, installing Node.js is necessary mainly to:
\begin{choices}
  \choice Replace the need for a web browser when viewing the app
  \CorrectChoice Run npm, the CLI, TypeScript compilation, and dev/test tooling locally
  \choice Compile CSS into TypeScript interfaces
  \choice Host the production database inside \texttt{node\_modules}
\end{choices}
\begin{solution}
  Install Node (LTS recommended), verify with \texttt{node -v} and \texttt{npm -v}, then \texttt{npm install -g @angular/cli} if your course uses global CLI. The app users load in Chrome still runs compiled JS, not Node.
\end{solution}

\newpage

\uplevel{\par\textbf{Part 6: Type annotations and basic types}\par}

\question[4]
Which declaration gives \texttt{name} a compile-time type of \texttt{string}?
\begin{choices}
  \choice \lstinline[style=tsinline]|var name = 42;|
  \CorrectChoice \lstinline[style=tsinline]|let name: string = 'Ada';|
  \choice \lstinline[style=tsinline]|const name: string;| with no initializer
  \choice \lstinline[style=tsinline]|name: string;| outside any class or function
\end{choices}
\begin{solution}
  Annotations follow the name: \lstinline[style=tsinline]|identifier: Type|. Type inference can omit the annotation when the initializer makes the type obvious (\lstinline[style=tsinline]|let n = 5| $\rightarrow$ \texttt{number}).
\end{solution}

\question[4]
The TypeScript type \lstinline[style=tsinline]|any|:
\begin{choices}
  \choice Forces strict checking on every operation
  \CorrectChoice Disables many compile-time checks for that value (escape hatch; avoid when possible)
  \choice Means the value must be \texttt{null}
  \choice Exists only at runtime, not in source files
\end{choices}
\begin{solution}
  Prefer precise types or \lstinline[style=tsinline]|unknown| (must narrow before use). Overusing \lstinline[style=tsinline]|any| removes much of TypeScript's safety benefit.
\end{solution}

\question[4]
A \textbf{tuple} type such as \lstinline[style=tsinline]|[string, number]| means:
\begin{choices}
  \choice An object with any number of string keys
  \CorrectChoice A fixed-length array with a known type at each index position
  \choice The same as \lstinline[style=tsinline]|Array<string>| only
  \choice A generic class constructor
\end{choices}
\begin{solution}
\begin{lstlisting}[style=tsexam]
let pair: [string, number] = ['score', 10];
let label = pair[0];  // string
let value = pair[1]; // number
\end{lstlisting}
  Tuples model ordered pairs such as \texttt{[name, age]} or API return tuples.
\end{solution}

\question[4]
The \lstinline[style=tsinline]|readonly| modifier on a property:
\begin{choices}
  \choice Allows reassignment after the object is created
  \CorrectChoice Prevents assignment to that property after initialization (compile-time check)
  \choice Converts the value to \lstinline[style=tsinline]|any| at runtime
  \choice Replaces the need for interfaces
\end{choices}
\begin{solution}
  \lstinline[style=tsinline]|readonly id: number| can be set in the constructor or at declaration, but not reassigned later. \lstinline[style=tsinline]|ReadonlyArray<T>| prevents mutating methods like \lstinline[style=tsinline]|push|.
\end{solution}

\question[4]
With \textbf{strict null checks} enabled, \lstinline[style=tsinline]|string| and \lstinline[style=tsinline]|null|:
\begin{choices}
  \choice Are always interchangeable
  \CorrectChoice Are distinct unless you use a union such as \lstinline[style=tsinline]|string | null|
  \choice Cannot appear in the same program
  \choice Apply only to generic classes
\end{choices}
\begin{solution}
  Strict mode catches \lstinline[style=tsinline]|null|/\lstinline[style=tsinline]|undefined| mistakes at compile time. Optional properties (\lstinline[style=tsinline]|name?: string|) include \lstinline[style=tsinline]|undefined| in their type.
\end{solution}

\newpage

\uplevel{\par\textbf{Part 7: Interfaces, classes, and enums}\par}

\question[4]
An \textbf{interface} in TypeScript is used to:
\begin{choices}
  \choice Compile directly to machine code without JavaScript
  \CorrectChoice Describe the shape of an object (property names and types) for compile-time checking
  \choice Replace \lstinline[style=tsinline]|function| declarations entirely
  \choice Run only inside Angular templates
\end{choices}
\begin{solution}
\begin{lstlisting}[style=tsexam]
interface User {
  id: number;
  name: string;
  email?: string;
}
function greet(u: User) {
  return u.name.toUpperCase();
}
\end{lstlisting}
  Interfaces are erased when compiling to JavaScript---they exist for the type checker only.
\end{solution}

\question[4]
Optional interface properties are marked with:
\begin{choices}
  \choice A leading \texttt{@} symbol
  \CorrectChoice \lstinline[style=tsinline]|?| after the property name
  \choice The keyword \lstinline[style=tsinline]|optional| before \lstinline[style=tsinline]|class|
  \choice Square brackets around the whole interface
\end{choices}
\begin{solution}
  \lstinline[style=tsinline]|email?: string| means callers may omit \texttt{email}; reads may yield \lstinline[style=tsinline]|undefined|. Combine with default values or null checks in implementations.
\end{solution}

\question[4]
A class \lstinline[style=tsinline]|implements| an interface when:
\begin{choices}
  \choice The class extends the interface at runtime
  \CorrectChoice The class agrees to provide every required property and method the interface declares
  \choice The interface inherits from the class automatically
  \choice The class is declared inside the interface file only
\end{choices}
\begin{solution}
\begin{lstlisting}[style=tsexam]
interface Printable { print(): void; }
class Report implements Printable {
  print() { console.log('report'); }
}
\end{lstlisting}
  \lstinline[style=tsinline]|implements| is checked by TypeScript; it does not copy implementation from the interface.
\end{solution}

\question[4]
In a TypeScript class, \lstinline[style=tsinline]|private| members:
\begin{choices}
  \choice Are visible on every subclass and external object by default
  \CorrectChoice Are accessible only inside the declaring class (TypeScript visibility check at compile time)
  \choice Cannot exist in the same class as \lstinline[style=tsinline]|public| members
  \choice Are the same as JavaScript private fields with \texttt{\#} only
\end{choices}
\begin{solution}
  \lstinline[style=tsinline]|public| (default), \lstinline[style=tsinline]|protected| (class + subclasses), and \lstinline[style=tsinline]|private| control access. ECMAScript \lstinline[style=tsinline]|#field| is a separate runtime private mechanism.
\end{solution}

\question[4]
A numeric \textbf{enum} such as:
\begin{lstlisting}[style=tsexam]
enum Direction { Up, Down, Left, Right }
\end{lstlisting}
\begin{choices}
  \choice Cannot be used as a type annotation
  \CorrectChoice Defines a set of named constants and generates a corresponding JavaScript object at compile time
  \choice Is identical to a TypeScript interface
  \choice Replaces \lstinline[style=tsinline]|keyof| on all objects
\end{choices}
\begin{solution}
  By default, members receive \texttt{0, 1, 2, \ldots}. You can assign explicit values. Use \lstinline[style=tsinline]|Direction.Up| for readable code instead of magic numbers.
\end{solution}

\question[4]
A \textbf{type alias} (\lstinline[style=tsinline]|type Point = { x: number; y: number }|) compared with an \textbf{interface}:
\begin{choices}
  \choice Can only describe primitives, never objects
  \CorrectChoice Can alias unions and primitives; interfaces are often preferred for object shapes that may be extended
  \choice Is checked at runtime by the browser
  \choice Must be declared in every \texttt{.html} file
\end{choices}
\begin{solution}
  Use \lstinline[style=tsinline]|type| for unions (\lstinline[style=tsinline]|string | number|), tuples, and mapped types. Use \lstinline[style=tsinline]|interface| for extendable object contracts (\lstinline[style=tsinline]|interface B extends A {}|).
\end{solution}

\uplevel{\par\textbf{Part 8: Unions, narrowing, and type guards}\par}

\question[4]
The union type \lstinline[style=tsinline]|string | number| means a value may be:
\begin{choices}
  \choice Both a string and a number at the same time always
  \CorrectChoice Either a \texttt{string} or a \texttt{number} (one or the other per value)
  \choice Only \lstinline[style=tsinline]|null|
  \choice A generic type parameter \texttt{T} only
\end{choices}
\begin{solution}
  Unions express alternatives. The compiler allows only members \textbf{common to all} union parts until you \emph{narrow} the type.
\end{solution}

\question[4]
Which statement \textbf{narrows} a union \lstinline[style=tsinline]|string | number| safely?
\begin{choices}
  \choice Call \lstinline[style=tsinline]|value.toFixed()| without checking
  \CorrectChoice \lstinline[style=tsinline]|if (typeof value === 'number') { value.toFixed(2); }|
  \choice Assign \lstinline[style=tsinline]|value = null| to widen the union
  \choice Cast with \lstinline[style=tsinline]|as any| in every branch
\end{choices}
\begin{solution}
  \textbf{typeof}, \textbf{instanceof}, \textbf{in}, and user-defined \textbf{type predicates} (\lstinline[style=tsinline]|value is Fish|) tell the checker which branch is active.
\end{solution}

\question[4]
A \textbf{type predicate} function looks like:
\begin{choices}
  \choice \lstinline[style=tsinline]|function isNum(x: any) { return true; }|
  \CorrectChoice \lstinline[style=tsinline]|function isFish(pet: Fish | Bird): pet is Fish { return (pet as Fish).swim !== undefined; }|
  \choice \lstinline[style=tsinline]|interface is Fish { }|
  \choice \lstinline[style=tsinline]|enum isFish { Yes }|
\end{choices}
\begin{solution}
  The return type \lstinline[style=tsinline]|pet is Fish| signals narrowing to callers. After \lstinline[style=tsinline]|if (isFish(p))|), TypeScript treats \lstinline[style=tsinline]|p| as \lstinline[style=tsinline]|Fish| in the true branch.
\end{solution}

\question[4]
The \lstinline[style=tsinline]|never| type represents:
\begin{choices}
  \choice Any possible value in the program
  \CorrectChoice Values that never occur (for example a function that always throws, or an impossible union branch)
  \choice The same as \lstinline[style=tsinline]|undefined| only
  \choice A deprecated alias for \lstinline[style=tsinline]|string|
\end{choices}
\begin{solution}
  Exhaustive \lstinline[style=tsinline]|switch| checks use \lstinline[style=tsinline]|never| in the \lstinline[style=tsinline]|default| arm to catch unhandled cases when a union grows later.
\end{solution}

\question[4]
The \lstinline[style=tsinline]|unknown| type differs from \lstinline[style=tsinline]|any| because \lstinline[style=tsinline]|unknown|:
\begin{choices}
  \choice Allows any operation without checks
  \CorrectChoice Requires narrowing or assertion before most operations (safer than \lstinline[style=tsinline]|any|)
  \choice Exists only inside enums
  \choice Cannot be returned from a function
\end{choices}
\begin{solution}
  Use \lstinline[style=tsinline]|unknown| for values from \lstinline[style=tsinline]|JSON.parse|, user input, or external APIs until validated.
\end{solution}

\question[4]
Literal types such as \lstinline[style=tsinline]|'success' | 'error'| are useful because they:
\begin{choices}
  \choice Force every string at runtime to be exactly six characters
  \CorrectChoice Restrict a value to specific string (or number) constants for safer APIs and exhaustiveness checking
  \choice Replace interfaces entirely
  \choice Compile to SQL \texttt{ENUM} columns automatically
\end{choices}
\begin{solution}
\begin{lstlisting}[style=tsexam]
type Status = 'loading' | 'ready' | 'error';
function setStatus(s: Status) { /* ... */ }
setStatus('ready');   // OK
setStatus('pending'); // compile error
\end{lstlisting}
\end{solution}

\newpage

\uplevel{\par\textbf{Part 9: Generics}\par}

\question[4]
A \textbf{generic class} declares type parameters:
\begin{choices}
  \choice In square brackets \texttt{[]} before the \texttt{class} keyword
  \CorrectChoice In angle brackets \texttt{<>} following the class name
  \choice Only on static members
  \choice In curly braces after the class body
\end{choices}
\begin{solution}
  Type parameters use angle brackets after the class name:
\begin{lstlisting}[style=tsexam]
class GenericNumber<NumType> {
  zeroValue: NumType;
  add: (x: NumType, y: NumType) => NumType;
}
let n = new GenericNumber<number>();
let s = new GenericNumber<string>();
\end{lstlisting}
  One parameter keeps fields and methods on the same type. Static members cannot use the class type parameter.
\end{solution}

\question[4]
To access \lstinline[style=tsinline]{arg.length} in a generic function, which constraint is used?
\begin{choices}
  \choice Unconstrained \texttt{<Type>} because every type has \texttt{length}
  \CorrectChoice \texttt{<Type extends Lengthwise>} where \texttt{Lengthwise} requires \texttt{length: number}
  \choice \texttt{any} only, with no interface
  \choice A SQL \texttt{JOIN}
\end{choices}
\begin{solution}
  Unconstrained \lstinline[style=tsinline]{<Type>} cannot access \lstinline[style=tsinline]{length}. Use a constraint:
\begin{lstlisting}[style=tsexam]
interface Lengthwise { length: number; }
function loggingIdentity<T extends Lengthwise>(arg: T): T {
  console.log(arg.length);
  return arg;
}
\end{lstlisting}
  \lstinline[style=tsinline]|loggingIdentity(3)| fails (a number has no \lstinline[style=tsinline]|length|). \lstinline[style=tsinline]|loggingIdentity({ length: 10, value: 3 })| is valid.
\end{solution}

\question[4]
\texttt{getProperty<Type, Key extends keyof Type>(obj: Type, key: Key)} ensures:
\begin{choices}
  \choice \texttt{key} may be any string
  \CorrectChoice \texttt{key} must be a property name that exists on \texttt{obj}
  \choice \texttt{obj} must be a number array
  \choice \texttt{keyof} is runtime-only
\end{choices}
\begin{solution}
\begin{lstlisting}[style=tsexam]
function getProperty<T, K extends keyof T>(obj: T, key: K) {
  return obj[key];
}
let x = { a: 1, b: 2, c: 3, d: 4 };
getProperty(x, "a");  // OK
getProperty(x, "m");  // compile error
\end{lstlisting}
\end{solution}

\question[4]
When creating generic \textbf{factory} functions in TypeScript, you refer to class types by:
\begin{choices}
  \choice Instance methods only
  \CorrectChoice Constructor signatures (for example \texttt{\{ new (): Type \}} or \texttt{new () => A})
  \choice Global \texttt{window} variables
  \choice SQL tables
\end{choices}
\begin{solution}
  A factory accepts a constructor, not an instance:
\begin{lstlisting}[style=tsexam]
function create<T>(c: { new(): T }): T {
  return new c();
}
function createInstance<A extends Animal>(c: new () => A): A {
  return new c();
}
\end{lstlisting}
  \lstinline[style=tsinline]{createInstance(Lion).keeper.nametag} type-checks because \lstinline[style=tsinline]{A} is inferred from the class passed in.
\end{solution}

\question[4]
\texttt{KeyValuePair<T, U>} with \texttt{setKeyValue(key: T, val: U)} illustrates that:
\begin{choices}
  \choice \texttt{T} and \texttt{U} must both be \texttt{number}
  \CorrectChoice One generic class can use different type pairs (for example \texttt{number}/\texttt{string} or \texttt{string}/\texttt{string}); a class may implement a generic interface and let the consumer bind types
  \choice Generic classes cannot implement interfaces
  \choice Only static fields may be generic
\end{choices}
\begin{solution}
\begin{lstlisting}[style=tsexam]
class KeyValuePair<T, U> {
  private key: T;
  private val: U;
  setKeyValue(key: T, val: U): void { /* body */ }
}
let kvp = new KeyValuePair<number, string>();
\end{lstlisting}
  Generics keep APIs type-safe while allowing different \lstinline[style=tsinline]{T}/\lstinline[style=tsinline]{U} pairs per use.
\end{solution}

\question[4]
Generic classes are only generic over their instance side rather than their static side.
\begin{choices}
  \CorrectChoice True
  \choice False
\end{choices}
\begin{solution}
  \textbf{True.} Type parameters on \lstinline[style=tsinline]{class Box<T>} apply to instance members only:
\begin{lstlisting}[style=tsexam]
class Box<T> {
  value: T;           // OK (instance)
  static bad: T;     // error: T not in scope on static side
}
\end{lstlisting}
\end{solution}

\question[4]
When creating factories in TypeScript using generics, it is necessary to refer to class types by their constructor functions.
\begin{choices}
  \CorrectChoice True
  \choice False
\end{choices}
\begin{solution}
  \textbf{True.} The factory parameter is a constructor signature; you pass the class and call \lstinline[style=tsinline]{new c()}:
\begin{lstlisting}[style=tsexam]
function create<T>(c: { new(): T }): T {
  return new c();
}
\end{lstlisting}
\end{solution}

\question[4]
When declaring a generic class, the generic type parameter is specified in square brackets after the name of the class.
\begin{choices}
  \choice True
  \CorrectChoice False
\end{choices}
\begin{solution}
  \textbf{False.} Generics use angle brackets:
\begin{lstlisting}[style=tsexam]
class KeyValuePair<T, U> { ... }   // correct
class Bad[T, U] { ... }            // wrong syntax
\end{lstlisting}
  Square brackets \lstinline[style=tsinline]{[]} are for arrays and indexing, not type parameters.
\end{solution}

\newpage

\uplevel{\par\textbf{Part 10: \texttt{tsconfig.json} basics}\par}

\question[4]
A \texttt{tsconfig.json} file tells the TypeScript compiler:
\begin{choices}
  \choice Which npm packages to download into \texttt{node\_modules}
  \CorrectChoice Which source files belong to the project and which compiler options to apply
  \choice How to render Angular templates in the browser
  \choice The CSS encapsulation mode for each component
\end{choices}
\begin{solution}
  Running \texttt{tsc} with no file arguments searches upward for \texttt{tsconfig.json} and builds the \textbf{project} described there. Single-file \texttt{tsc file.ts} bypasses project settings unless you pass \texttt{-p}.
\end{solution}

\question[4]
Top-level \texttt{include} and \texttt{exclude} in \texttt{tsconfig.json} control:
\begin{choices}
  \choice Which npm scripts appear in \texttt{package.json}
  \CorrectChoice Which files are part of the compilation (globs such as \texttt{src/**/*} vs \texttt{node\_modules})
  \choice Runtime routing paths in Angular
  \choice Karma browser selection only
\end{choices}
\begin{solution}
\begin{lstlisting}[style=jsonexam]
{
  "include": ["src/**/*"],
  "exclude": ["node_modules", "dist", "**/*.spec.ts"]
}
\end{lstlisting}
  If \texttt{include} is omitted, \texttt{tsc} defaults to all \texttt{.ts} files in the folder and subfolders except \texttt{node\_modules}, \texttt{bower\_components}, and \texttt{jspm\_packages}.
\end{solution}

\question[4]
The \texttt{files} array in \texttt{tsconfig.json}:
\begin{choices}
  \choice Lists every file in \texttt{node\_modules} recursively
  \CorrectChoice Explicitly names entry \texttt{.ts} files to compile when you want a fixed list instead of globs
  \choice Replaces \texttt{compilerOptions} entirely
  \choice Must contain only \texttt{.css} paths
\end{choices}
\begin{solution}
  \texttt{include}/\texttt{exclude} globs are more common in apps. \texttt{files} is useful for small libraries with a handful of entry points. You cannot mix \texttt{files} with \texttt{include} in all toolchain versions---follow one style per project.
\end{solution}

\question[4]
\texttt{"extends": "./tsconfig.base.json"} in a \texttt{tsconfig.json}:
\begin{choices}
  \choice Imports JavaScript at runtime
  \CorrectChoice Inherits settings from another config file and merges overrides (Angular workspaces often use a base + app configs)
  \choice Forces \texttt{target} to \texttt{ES5} only
  \choice Disables \texttt{strict} permanently
\end{choices}
\begin{solution}
  Child configs override parent keys. Monorepos may have a root base and per-package \texttt{tsconfig.app.json}, \texttt{tsconfig.spec.json}, etc.
\end{solution}

\question[4]
\texttt{compilerOptions} in \texttt{tsconfig.json} should be placed:
\begin{choices}
  \choice Inside \texttt{package.json} \texttt{scripts} only
  \CorrectChoice Inside the JSON object, sibling to \texttt{include}/\texttt{exclude} (not outside the root)
  \choice In each \texttt{.html} template
  \choice In \texttt{angular.json} \texttt{styles} array only
\end{choices}
\begin{solution}
  Typical layout: top-level \texttt{compilerOptions}, optional \texttt{include}/\texttt{exclude}/\texttt{files}, optional \texttt{extends}. All flags that control emit and type checking live under \texttt{compilerOptions}.
\end{solution}

\question[4]
Angular CLI projects usually compile TypeScript using:
\begin{choices}
  \choice A hand-written \texttt{tsc file.ts} per component with no config
  \CorrectChoice Tooling driven by \texttt{tsconfig.app.json} / \texttt{tsconfig.spec.json} (or equivalent) referenced from the build
  \choice Only Babel with no \texttt{tsconfig}
  \choice \texttt{ng serve} without any TypeScript configuration files
\end{choices}
\begin{solution}
  You still benefit from understanding \texttt{tsconfig}: IDE diagnostics, library builds, and custom packages outside the CLI all rely on the same options.
\end{solution}

\uplevel{\par\textbf{Part 11: Compiler options}\par}

\question[4]
The \texttt{outDir} compiler option specifies:
\begin{choices}
  \choice The folder where \texttt{package.json} must live
  \CorrectChoice Where emitted JavaScript files are written (for example \texttt{dist/})
  \choice The DOM mount point for Angular
  \choice The Karma port
\end{choices}
\begin{solution}
  Pair with \texttt{rootDir} (default input root, often \texttt{src}) so \texttt{src/app.ts} maps to \texttt{dist/app.js} instead of flattening unexpectedly.
\end{solution}

\question[4]
The \texttt{module} compiler option controls:
\begin{choices}
  \choice CSS module scoping in components
  \CorrectChoice The module format of generated \texttt{import}/\texttt{export} code (for example \texttt{commonjs}, \texttt{esnext}, \texttt{nodenext})
  \choice The Angular \texttt{NgModule} class name
  \choice Whether \texttt{strict} is enabled
\end{choices}
\begin{solution}
  Match \texttt{module} to your runtime/bundler: Node historically used \texttt{commonjs}; modern bundlers often use \texttt{esnext} or \texttt{es2015}+. Mismatch causes \texttt{require} vs \texttt{import} errors at run time.
\end{solution}

\question[4]
Setting \texttt{"declaration": true} makes \texttt{tsc} emit:
\begin{choices}
  \choice HTML templates for each class
  \CorrectChoice \texttt{.d.ts} declaration files describing types for consumers of a library
  \choice Source maps only
  \choice SQL schema files
\end{choices}
\begin{solution}
  Published npm libraries often ship \texttt{.d.ts} so TypeScript users get autocomplete without source. Apps targeting browsers may omit declarations.
\end{solution}

\question[4]
The \texttt{lib} option (for example \texttt{["ES2020", "DOM"]}) tells TypeScript:
\begin{choices}
  \choice Which npm packages to install automatically
  \CorrectChoice Which built-in API typings are available (\texttt{Array}, \texttt{Promise}, \texttt{document}, etc.)
  \choice The ECMAScript \texttt{target} emit version only
  \choice The path to \texttt{node\_modules}
\end{choices}
\begin{solution}
  \texttt{target} controls emitted JS syntax; \texttt{lib} controls which \emph{type definitions} for globals you may use. A Node project might use \texttt{["ES2020"]} without \texttt{DOM}; a browser app adds \texttt{DOM}.
\end{solution}

\question[4]
\texttt{esModuleInterop: true} helps when:
\begin{choices}
  \choice Disabling all type checks
  \CorrectChoice Importing CommonJS modules with default imports in a more ergonomic way (\texttt{import fs from 'fs'} style)
  \choice Removing the need for \texttt{target}
  \choice Compiling CSS to TypeScript
\end{choices}
\begin{solution}
  Many configs also set \texttt{allowSyntheticDefaultImports} for editor compatibility. These options reduce friction between module systems, not replace learning how \texttt{import} works.
\end{solution}

\question[4]
\texttt{sourceMap: true} causes \texttt{tsc} to emit:
\begin{choices}
  \choice Only \texttt{.d.ts} files
  \CorrectChoice \texttt{.js.map} files that link compiled JavaScript back to TypeScript lines for debugging
  \choice Angular route configs
  \choice Extra \texttt{.html} copies of templates
\end{choices}
\begin{solution}
  Browser devtools can show original \texttt{.ts} line numbers when maps are present. Production builds sometimes omit maps or ship them separately for security/size.
\end{solution}

\uplevel{\par\textbf{Part 12: \texttt{strict} mode and type safety}\par}

\question[4]
Setting \texttt{"strict": true} in \texttt{compilerOptions}:
\begin{choices}
  \choice Turns off all errors so builds always succeed
  \CorrectChoice Enables a family of stricter checks (including \texttt{strictNullChecks} and \texttt{noImplicitAny}) recommended for new projects
  \choice Forces \texttt{target} to \texttt{ES3}
  \choice Prevents emitting JavaScript files
\end{choices}
\begin{solution}
  \texttt{strict} is a shortcut flag. You can toggle individual strict sub-options when migrating legacy code, but greenfield TypeScript (and Angular CLI defaults) favor strict on.
\end{solution}

\question[4]
\texttt{strictNullChecks: true} means:
\begin{choices}
  \choice \lstinline[style=tsinline]|null| and \lstinline[style=tsinline]|undefined| may be assigned to any type freely
  \CorrectChoice \lstinline[style=tsinline]|null|/\lstinline[style=tsinline]|undefined| are not assignable to other types unless the type explicitly allows them (for example \lstinline[style=tsinline]|string | null|)
  \choice All functions must return \lstinline[style=tsinline]|never|
  \choice Interfaces are erased at runtime
\end{choices}
\begin{solution}
  This catches \texttt{Cannot read property 'x' of null} class bugs earlier. Use optional chaining (\lstinline[style=tsinline]|obj?.field|) and nullish coalescing (\lstinline[style=tsinline]|value ?? default|) with strict code.
\end{solution}

\question[4]
With \texttt{noImplicitAny: true}, a parameter with no type annotation and no inferred type:
\begin{choices}
  \choice Defaults to \lstinline[style=tsinline]|never|
  \CorrectChoice Causes a compile error instead of silently becoming \lstinline[style=tsinline]|any|
  \choice Is deleted from the emit file
  \choice Becomes \lstinline[style=tsinline]|string| automatically
\end{choices}
\begin{solution}
  Explicit \lstinline[style=tsinline]|any| is still allowed when you intend it. The flag stops \emph{accidental} \lstinline[style=tsinline]|any| from untyped parameters and untyped variables.
\end{solution}

\question[4]
\texttt{strictFunctionTypes: true} tightens checking of:
\begin{choices}
  \choice CSS class names in templates
  \CorrectChoice Function parameter types (parameters are checked more strictly when assigning functions)
  \choice JSON \texttt{stringify} output
  \choice npm semver strings only
\end{choices}
\begin{solution}
  Part of the \texttt{strict} bundle. It reduces unsound callbacks when passing functions as arguments---common in event handlers and RxJS operators.
\end{solution}

\question[4]
Turning \texttt{strict} off in a large legacy codebase:
\begin{choices}
  \choice Automatically fixes all runtime bugs
  \CorrectChoice May allow unsound code to compile but removes many helpful compile-time guards
  \choice Deletes \texttt{node\_modules}
  \choice Forces \texttt{target ES5} only
\end{choices}
\begin{solution}
  Migration strategy: enable \texttt{strict} in new folders, use \texttt{// @ts-expect-error} sparingly with tickets, or enable sub-flags one at a time (\texttt{strictNullChecks} first is common).
\end{solution}

\question[4]
\texttt{noUnusedLocals} and \texttt{noUnusedParameters} (when enabled) report:
\begin{choices}
  \choice Only runtime exceptions in the browser
  \CorrectChoice Variables or parameters that are declared but never read (as compile errors or warnings depending on setup)
  \choice Missing \texttt{package.json} fields
  \choice Invalid HTML selectors
\end{choices}
\begin{solution}
  Not part of \texttt{strict} by default but useful for hygiene. Prefix unused parameters with \lstinline[style=tsinline]|_| (\lstinline[style=tsinline]|_event|) if a signature must match an interface.
\end{solution}

\newpage

\uplevel{\par\textbf{Part 13: \texttt{target} and emitted JavaScript}\par}

\question[4]
The \texttt{target} compiler option (for example \texttt{ES2017}) sets:
\begin{choices}
  \choice Which browsers may download \texttt{package.json}
  \CorrectChoice The ECMAScript \textbf{version of JavaScript syntax} that \texttt{tsc} emits
  \choice The folder name for unit tests only
  \choice Whether Angular uses Ivy
\end{choices}
\begin{solution}
  Higher targets emit modern syntax (\texttt{async}/\texttt{await}, classes, arrow functions) with less downleveling. Older targets produce more helper code for legacy browsers.
\end{solution}

\question[4]
If \texttt{target} is \texttt{ES5}, \texttt{tsc} may transform modern features such as:
\begin{choices}
  \choice HTML templates into SQL
  \CorrectChoice Classes and async functions into older-compatible JavaScript patterns
  \choice All \texttt{import} statements into global variables only without helpers
  \choice Type annotations into runtime objects
\end{choices}
\begin{solution}
  Downleveling increases bundle size. Many Angular apps target evergreen browsers and use a modern \texttt{target} plus a bundler; support for IE11 required much older targets and polyfills.
\end{solution}

\question[4]
\texttt{target} and \texttt{lib} differ because \texttt{target}:
\begin{choices}
  \CorrectChoice Controls emitted JS language level, while \texttt{lib} controls which type definitions for standard APIs are available during checking
  \choice Is the same key with two names
  \choice Only applies to \texttt{.css} files
  \choice Replaces \texttt{moduleResolution}
\end{choices}
\begin{solution}
  You can \texttt{target ES5} but include \texttt{lib ES2020} only if the runtime provides polyfills---usually keep them aligned with your minimum browser/Node version.
\end{solution}

\question[4]
Choosing \texttt{ES2020} as \texttt{target} is reasonable when:
\begin{choices}
  \choice You must support Internet Explorer 8 exclusively
  \CorrectChoice Your deployment environment (modern browsers or current Node LTS) supports that language level natively
  \choice You want to disable type checking
  \choice You never run \texttt{tsc}
\end{choices}
\begin{solution}
  Check browser support matrices and Node documentation. Angular CLI sets sensible defaults; custom libraries may target older \texttt{target} for wider npm consumers.
\end{solution}

\question[4]
The \texttt{useDefineForClassFields} option relates to:
\begin{choices}
  \choice Karma reporters only
  \CorrectChoice How class fields are emitted to match ECMAScript class field semantics vs older TypeScript behavior
  \choice Router lazy loading
  \choice \texttt{package-lock.json} format
\end{choices}
\begin{solution}
  Modern \texttt{target} + this flag align field initialization with the JavaScript spec. Misconfiguration can change runtime initialization order in classes---relevant when mixing TS classes and inheritance.
\end{solution}

\question[4]
After changing \texttt{target} from \texttt{ES5} to \texttt{ES2018}, you should:
\begin{choices}
  \choice Delete \texttt{tsconfig.json}
  \CorrectChoice Rebuild (\texttt{tsc} or \texttt{ng build}) and retest in your minimum supported environment
  \choice Remove all \texttt{strict} flags automatically
  \choice Stop using \texttt{module} settings
\end{choices}
\begin{solution}
  Emit changes affect syntax (for example native \texttt{async}/\texttt{await}). Bundlers may still transpile further via Babel if configured---know both \texttt{tsc} \texttt{target} and downstream tooling.
\end{solution}

\uplevel{\par\textbf{Part 14: The \texttt{tsc} compiler in depth}\par}

\question[4]
Running \texttt{tsc} in a folder that contains \texttt{tsconfig.json} with no extra arguments:
\begin{choices}
  \choice Compiles only \texttt{main.ts} and ignores the config
  \CorrectChoice Builds the whole project described by the config (all included files)
  \choice Starts the Angular dev server on port \texttt{4200}
  \choice Installs npm dependencies
\end{choices}
\begin{solution}
  Contrast: \texttt{tsc src/app.ts} compiles one file with default options unless you pass \texttt{--project}. Prefer project builds for real apps.
\end{solution}

\question[4]
The flag \texttt{tsc --noEmit}:
\begin{choices}
  \choice Deletes all \texttt{.ts} sources
  \CorrectChoice Type-checks the project without writing JavaScript output files
  \choice Disables \texttt{strict} permanently
  \choice Emits only \texttt{.css} files
\end{choices}
\begin{solution}
  CI pipelines and editors use typecheck-only runs. Fast feedback without updating \texttt{dist/}. Pair with \texttt{--pretty} for readable errors.
\end{solution}

\question[4]
To compile using \texttt{tsconfig.app.json} instead of the default file name, use:
\begin{choices}
  \choice \texttt{tsc --init tsconfig.app.json}
  \CorrectChoice \texttt{tsc -p tsconfig.app.json} (or \texttt{--project})
  \choice \texttt{tsc --target tsconfig.app.json}
  \choice \texttt{npm install tsconfig.app.json}
\end{choices}
\begin{solution}
  Angular separates app, spec, and library configs. Each \texttt{-p} path is one TypeScript \textbf{project} with its own \texttt{include} and \texttt{compilerOptions}.
\end{solution}

\question[4]
\texttt{tsc --showConfig} prints:
\begin{choices}
  \choice The contents of \texttt{node\_modules}
  \CorrectChoice The fully merged configuration after \texttt{extends} and defaults are applied
  \choice Only the TypeScript version installed globally
  \choice Karma test results
\end{choices}
\begin{solution}
  Useful when debugging why a flag seems ignored---you see the effective \texttt{compilerOptions} object.
\end{solution}

\question[4]
If \texttt{tsc} reports type errors during a project build, you should generally:
\begin{choices}
  \choice Ignore them because JavaScript will fix types at runtime
  \CorrectChoice Treat the build as failed (non-zero exit code), fix the errors, then emit---do not ship unchecked code silently
  \choice Delete \texttt{tsconfig.json} and retry
  \choice Convert all errors to CSS warnings only
\end{choices}
\begin{solution}
  Use \texttt{tsc --noEmit} in CI for typecheck-only gates. Tools that skip checking (\texttt{transpileOnly}) trade speed for safety---know which step enforces types in your pipeline.
\end{solution}

\question[4]
\texttt{tsc --version} (or \texttt{npx tsc -v}) shows:
\begin{choices}
  \choice The Node.js patch level only
  \CorrectChoice The installed TypeScript compiler package version (should match the project's \texttt{devDependency} in teams)
  \choice The Angular CLI version always
  \choice The browser user agent
\end{choices}
\begin{solution}
  Mismatched global vs local \texttt{tsc} causes ``works on my machine'' issues. Scripts should call \texttt{npx tsc} or \texttt{./node\_modules/.bin/tsc} so everyone uses the same compiler version pinned in \texttt{package.json}.
\end{solution}

\question[4]
Building after \texttt{tsc -w} (\texttt{--watch}) is suited for:
\begin{choices}
  \choice Production deployment with no further bundling ever
  \CorrectChoice Local development when you want JavaScript output refreshed on every save
  \choice Running end-to-end browser tests only
  \choice Generating \texttt{package.json} from HTML
\end{choices}
\begin{solution}
  Watch mode re-runs the compiler on file changes. Angular CLI's \texttt{ng serve} integrates similar watching with webpack/esbuild; standalone libraries often use \texttt{tsc -w} plus a test runner.
\end{solution}

\newpage

\uplevel{\par\textbf{Part 15: Basic types (primitives and variables)}\par}

\question[4]
Which TypeScript types are considered \textbf{primitive} (simple) types?
\begin{choices}
  \choice Only \texttt{interface} and \texttt{class}
  \CorrectChoice \lstinline[style=tsinline]|string|, \lstinline[style=tsinline]|number|, \lstinline[style=tsinline]|boolean|, \lstinline[style=tsinline]|bigint|, \lstinline[style=tsinline]|symbol|, \lstinline[style=tsinline]|null|, and \lstinline[style=tsinline]|undefined|
  \choice Only \lstinline[style=tsinline]|Array| and \lstinline[style=tsinline]|Object|
  \choice Only Angular \lstinline[style=tsinline]|Component|
\end{choices}
\begin{solution}
  Primitives are immutable values, not object instances (except boxing at runtime). Most variables in apps use \lstinline[style=tsinline]|string|, \lstinline[style=tsinline]|number|, and \lstinline[style=tsinline]|boolean| daily.
\end{solution}

\question[4]
In TypeScript, \lstinline[style=tsinline]|let active: boolean = true;| declares:
\begin{choices}
  \choice A string that stores \texttt{"true"}
  \CorrectChoice A variable that may be reassigned and must hold only \texttt{true} or \texttt{false}
  \choice A constant that can never change
  \choice A numeric enum member
\end{choices}
\begin{solution}
  \lstinline[style=tsinline]|boolean| is for logic flags. \lstinline[style=tsinline]|const| prevents rebinding the variable; \lstinline[style=tsinline]|let| allows reassignment to another boolean value.
\end{solution}

\question[4]
The TypeScript \lstinline[style=tsinline]|number| type includes:
\begin{choices}
  \choice Only integers---decimals are a separate \texttt{decimal} type always
  \CorrectChoice Integers and floats (all IEEE-754 values), including \lstinline[style=tsinline]|NaN| and \lstinline[style=tsinline]|Infinity|
  \choice Only positive whole numbers
  \choice Only values that fit in 8 bits
\end{choices}
\begin{solution}
  Unlike some languages, there is no separate \texttt{int}/\texttt{float} keyword. For very large integers, use \lstinline[style=tsinline]|bigint| (\lstinline[style=tsinline]|100n| syntax).
\end{solution}

\question[4]
Type \textbf{inference} on a variable means:
\begin{choices}
  \choice The browser guesses types at runtime
  \CorrectChoice TypeScript deduces the type from the initializer when you omit an annotation (\lstinline[style=tsinline]|let count = 0| $\rightarrow$ \texttt{number})
  \choice You must write \lstinline[style=tsinline]|any| on every line
  \choice Only interfaces can be inferred
\end{choices}
\begin{solution}
\begin{lstlisting}[style=tsexam]
let title = 'Angular';     // inferred string
let count: number = 0;     // explicit number
const ids = [1, 2, 3];     // inferred number[]
\end{lstlisting}
  \lstinline[style=tsinline]|const| infers narrow literal types when possible (\lstinline[style=tsinline]|const x = 1| is literal \lstinline[style=tsinline]|1|, not general \lstinline[style=tsinline]|number|).
\end{solution}

\question[4]
Which array type annotation is idiomatic for an array of strings?
\begin{choices}
  \choice \lstinline[style=tsinline]|Array<strings>|
  \CorrectChoice \lstinline[style=tsinline]|string[]| or \lstinline[style=tsinline]|Array<string>|
  \choice \lstinline[style=tsinline]|string{}|
  \choice \lstinline[style=tsinline]|[string]|
\end{choices}
\begin{solution}
  \lstinline[style=tsinline]|string[]| is common in hand-written code. \lstinline[style=tsinline]|Array<string>| is equivalent. \lstinline[style=tsinline]|[string]| is a one-tuple, not an arbitrary-length array.
\end{solution}

\question[4]
Reassigning \lstinline[style=tsinline]|let value: string = 'hi';| later with \lstinline[style=tsinline]|value = 42;|:
\begin{choices}
  \choice Is always allowed because numbers coerce to strings
  \CorrectChoice Produces a type error---\lstinline[style=tsinline]|42| is not assignable to \lstinline[style=tsinline]|string|
  \choice Changes the type of \lstinline[style=tsinline]|value| to \lstinline[style=tsinline]|any| automatically
  \choice Requires \texttt{strict} to be off and on at the same time
\end{choices}
\begin{solution}
  Variable annotations fix the allowed set of values for that binding. Use a union (\lstinline[style=tsinline]|string | number|) if both kinds are intentional.
\end{solution}

\uplevel{\par\textbf{Part 16: Special types}\par}

\question[4]
The \lstinline[style=tsinline]|void| type is most often used for:
\begin{choices}
  \choice Variables that store \texttt{null} only
  \CorrectChoice Functions that return no meaningful value (\lstinline[style=tsinline]|return;| or implicit fall-through)
  \choice JSON keys that are missing
  \choice Generic type parameters on classes
\end{choices}
\begin{solution}
\begin{lstlisting}[style=tsexam]
function logMessage(msg: string): void {
  console.log(msg);
}
\end{lstlisting}
  \lstinline[style=tsinline]|void| means ignore the return value. Do not confuse with \lstinline[style=tsinline]|undefined| (a value a variable can hold).
\end{solution}

\question[4]
In TypeScript, \lstinline[style=tsinline]|null| and \lstinline[style=tsinline]|undefined|:
\begin{choices}
  \choice Are identical types with the same name
  \CorrectChoice Are distinct primitive types (\lstinline[style=tsinline]|null| is an intentional absence of object; \lstinline[style=tsinline]|undefined| often means not assigned or missing optional value)
  \choice Exist only after compilation to JavaScript enums
  \choice Cannot appear in union types
\end{choices}
\begin{solution}
  With \texttt{strictNullChecks}, assign them only when the type allows (\lstinline[style=tsinline]|string | null|). Optional parameters default to \lstinline[style=tsinline]|undefined| when omitted.
\end{solution}

\question[4]
A function declared to return \lstinline[style=tsinline]|never|:
\begin{choices}
  \choice May return any value
  \CorrectChoice Never returns normally (always throws, infinite loop, or exhausts all paths in a checker proof)
  \choice Is the same as \lstinline[style=tsinline]|void|
  \choice Must return \lstinline[style=tsinline]|null|
\end{choices}
\begin{solution}
  Use \lstinline[style=tsinline]|never| for exhaustive \lstinline[style=tsinline]|switch| defaults and error helpers. The compiler uses it to prove code paths are unreachable.
\end{solution}

\question[4]
Assigning a value of type \lstinline[style=tsinline]|any| to a variable typed \lstinline[style=tsinline]|string|:
\begin{choices}
  \choice Always fails at compile time
  \CorrectChoice Is allowed by the compiler (and may hide bugs) because \lstinline[style=tsinline]|any| disables checking on that value
  \choice Converts the value to \lstinline[style=tsinline]|never| automatically
  \choice Requires \lstinline[style=tsinline]|readonly|
\end{choices}
\begin{solution}
  Treat \lstinline[style=tsinline]|any| as an escape hatch from migration or untyped libraries. Prefer \lstinline[style=tsinline]|unknown| plus narrowing for safer external data.
\end{solution}

\question[4]
A variable of type \lstinline[style=tsinline]|unknown|:
\begin{choices}
  \choice May call any method without checks
  \CorrectChoice Cannot call methods or access properties until narrowed (for example with \lstinline[style=tsinline]|typeof| or a type guard)
  \choice Is identical to \lstinline[style=tsinline]|string|
  \choice Compiles to a special runtime class
\end{choices}
\begin{solution}
\begin{lstlisting}[style=tsexam]
function printLength(x: unknown) {
  if (typeof x === 'string') {
    console.log(x.length);
  }
}
\end{lstlisting}
\end{solution}

\question[4]
A \textbf{string literal type} \lstinline[style=tsinline]|'left' | 'right'| differs from general \lstinline[style=tsinline]|string| because:
\begin{choices}
  \choice It allows any string at compile time
  \CorrectChoice Only those exact string values are assignable unless widened
  \choice It stores binary data only
  \choice It replaces \lstinline[style=tsinline]|enum| at runtime with objects always
\end{choices}
\begin{solution}
  Literal types model fixed domains (alignment, HTTP methods, status codes). They pair with unions and discriminated unions for precise APIs.
\end{solution}

\uplevel{\par\textbf{Part 17: Object types}\par}

\question[4]
An inline \textbf{object type} in TypeScript looks like:
\begin{choices}
  \choice \lstinline[style=tsinline]|object = string|
  \CorrectChoice \lstinline[style=tsinline]|{ id: number; name: string }|
  \choice \lstinline[style=tsinline]|<Object>|
  \choice \lstinline[style=tsinline]|Object<string>|
\end{choices}
\begin{solution}
\begin{lstlisting}[style=tsexam]
let user: { id: number; name: string } = { id: 1, name: 'Ada' };
function printId(u: { id: number }) {
  console.log(u.id);
}
\end{lstlisting}
  Reusable shapes usually move to an \lstinline[style=tsinline]|interface| or \lstinline[style=tsinline]|type| alias.
\end{solution}

\question[4]
The type \lstinline[style=tsinline]|Record<string, number>| describes:
\begin{choices}
  \choice A tuple of two numbers
  \CorrectChoice An object with string keys and \lstinline[style=tsinline]|number| values (dictionary/map shape)
  \choice An array of strings only
  \choice A function returning \lstinline[style=tsinline]|void|
\end{choices}
\begin{solution}
  \lstinline[style=tsinline]|Record<Keys, Type>| is a built-in utility for string-keyed maps. An \textbf{index signature} \lstinline[style=tsinline]|{ [key: string]: number }| is similar but can mix with named properties under stricter rules.
\end{solution}

\question[4]
An index signature \lstinline[style=tsinline]|{ [prop: string]: boolean }| allows:
\begin{choices}
  \choice Only the property named \texttt{prop}
  \CorrectChoice Any string-keyed property with boolean values (plus other declared members if specified)
  \choice Only numeric array indices
  \choice SQL-style joins between tables
\end{choices}
\begin{solution}
  Useful for dynamic dictionaries (\texttt{flags['darkMode']}). Prefer \lstinline[style=tsinline]|Record| or mapped types when keys are known unions (\lstinline[style=tsinline]|'a' | 'b'|).
\end{solution}

\question[4]
Marking a property \lstinline[style=tsinline]|readonly id: number| on an object type:
\begin{choices}
  \choice Allows \lstinline[style=tsinline]|obj.id = 2| after creation
  \CorrectChoice Prevents assigning to \lstinline[style=tsinline]|id| after the object is created (compile-time)
  \choice Deletes the property at runtime
  \choice Makes \lstinline[style=tsinline]|id| optional
\end{choices}
\begin{solution}
  \lstinline[style=tsinline]|Readonly<T>| utility makes all properties readonly. Deep immutability requires nested \lstinline[style=tsinline]|readonly| or libraries/immutable patterns.
\end{solution}

\question[4]
Optional object properties use:
\begin{choices}
  \choice A leading \texttt{@} on the key
  \CorrectChoice \lstinline[style=tsinline]|?| after the name (\lstinline[style=tsinline]|email?: string|)
  \choice The keyword \lstinline[style=tsinline]|optional| inside \lstinline[style=tsinline]|class| only
  \choice Square brackets around the entire object
\end{choices}
\begin{solution}
  Callers may omit optional keys. Access yields \lstinline[style=tsinline]|string | undefined| under strict null checks---guard before use or provide defaults.
\end{solution}

\question[4]
For an object shape that other teams may \textbf{extend} in declaration merging, prefer:
\begin{choices}
  \choice A union of primitives
  \CorrectChoice An \lstinline[style=tsinline]|interface| (supports \lstinline[style=tsinline]|interface B extends A {}| and merging in some cases)
  \choice Only \lstinline[style=tsinline]|number[]|
  \choice \lstinline[style=tsinline]|never| only
\end{choices}
\begin{solution}
  \lstinline[style=tsinline]|type| aliases excel at unions/intersections; \lstinline[style=tsinline]|interface| excels at object extension. Both describe object shapes; pick based on composition needs.
\end{solution}

\newpage

\uplevel{\par\textbf{Part 18: Union types (composition and narrowing)}\par}

\question[4]
The union \lstinline[style=tsinline]|string | number| means the value at runtime is:
\begin{choices}
  \choice Always both types simultaneously
  \CorrectChoice Either a string \textbf{or} a number in a single value
  \choice Always \lstinline[style=tsinline]|null|
  \choice Only an object with two fields
\end{choices}
\begin{solution}
  Unions widen what you can \textbf{assign in}. You must \textbf{narrow} before using members unique to one branch (see also Part~8).
\end{solution}

\question[4]
An \textbf{intersection} type \lstinline[style=tsinline]|A & B| means a value must:
\begin{choices}
  \choice Satisfy \textbf{either} \lstinline[style=tsinline]|A| or \lstinline[style=tsinline]|B| only
  \CorrectChoice Satisfy \textbf{both} \lstinline[style=tsinline]|A| and \lstinline[style=tsinline]|B| at the same time
  \choice Be \lstinline[style=tsinline]|never|
  \choice Ignore all properties from \lstinline[style=tsinline]|B|
\end{choices}
\begin{solution}
  Intersections merge capabilities: \lstinline[style=tsinline]|Named & Aged| has both name and age fields. Conflicting property types in intersections collapse to \lstinline[style=tsinline]|never| on that property.
\end{solution}

\question[4]
A \textbf{discriminated union} typically includes:
\begin{choices}
  \choice Only primitive strings with no shared field
  \CorrectChoice A common literal-typed field (for example \lstinline[style=tsinline]|kind: 'circle'| vs \lstinline[style=tsinline]|kind: 'square'|) so \lstinline[style=tsinline]|switch| can narrow safely
  \choice Only \lstinline[style=tsinline]|any| members
  \choice A \texttt{bootstrap} array
\end{choices}
\begin{solution}
\begin{lstlisting}[style=tsexam]
type Shape =
  | { kind: 'circle'; radius: number }
  | { kind: 'square'; side: number };

function area(s: Shape) {
  switch (s.kind) {
    case 'circle': return Math.PI * s.radius ** 2;
    case 'square': return s.side * s.side;
  }
}
\end{lstlisting}
\end{solution}

\question[4]
To narrow \lstinline[style=tsinline]|Fish | Bird| when both are interfaces, a reliable check is:
\begin{choices}
  \choice \lstinline[style=tsinline]|typeof pet === 'Fish'|
  \CorrectChoice \lstinline[style=tsinline]|'swim' in pet| or a user-defined type predicate (\lstinline[style=tsinline]|pet is Fish|)
  \choice \lstinline[style=tsinline]|pet == null| only
  \choice \lstinline[style=tsinline]|JSON.stringify(pet)| only
\end{choices}
\begin{solution}
  \lstinline[style=tsinline]|typeof| works for primitives. For objects, use \lstinline[style=tsinline]|in|, \lstinline[style=tsinline]|instanceof| (classes), or discriminant fields. Part~8 covers type predicates in more detail.
\end{solution}

\question[4]
Assigning \lstinline[style=tsinline]|'pending'| to a variable of type \lstinline[style=tsinline]|'success' | 'error'|:
\begin{choices}
  \choice Always succeeds because all are strings
  \CorrectChoice Fails type checking---\lstinline[style=tsinline]|'pending'| is not in the allowed literal union
  \choice Widens the variable to \lstinline[style=tsinline]|any| silently always
  \choice Requires \texttt{target ES5}
\end{choices}
\begin{solution}
  Literal unions are stricter than \lstinline[style=tsinline]|string|. They catch typos in status codes and event names at compile time.
\end{solution}

\question[4]
Exhaustiveness checking in a \lstinline[style=tsinline]|switch| on a union uses \lstinline[style=tsinline]|never| to:
\begin{choices}
  \choice Skip the \lstinline[style=tsinline]|default| branch always
  \CorrectChoice Prove every union member was handled---if a new variant is added, the \lstinline[style=tsinline]|default| arm errors when assigned to \lstinline[style=tsinline]|never|
  \choice Convert unions to intersections
  \choice Emit \texttt{.d.ts} files only
\end{choices}
\begin{solution}
\begin{lstlisting}[style=tsexam]
function assertNever(x: never): never {
  throw new Error('Unexpected value: ' + x);
}
\end{lstlisting}
  Pattern keeps \lstinline[style=tsinline]|switch| statements maintainable as unions grow.
\end{solution}

\newpage

\uplevel{\par\textbf{Part 19: Type aliases}\par}

\question[4]
A \textbf{type alias} is declared with:
\begin{choices}
  \choice \lstinline[style=tsinline]|interface Name = ...|
  \CorrectChoice \lstinline[style=tsinline]|type Name = ...;|
  \choice \lstinline[style=tsinline]|class Name = ...|
  \choice \lstinline[style=tsinline]|enum Name = ...| for every alias
\end{choices}
\begin{solution}
  Aliases create a new name for an existing type expression. They do not create runtime values (unlike \lstinline[style=tsinline]|enum| in many cases).
\end{solution}

\question[4]
Which definition is a valid \textbf{type alias} for a union?
\begin{choices}
  \choice \lstinline[style=tsinline]|interface Id = string | number;|
  \CorrectChoice \lstinline[style=tsinline]|type Id = string | number;|
  \choice \lstinline[style=tsinline]|type Id: string | number;|
  \choice \lstinline[style=tsinline]|union Id { string, number }|
\end{choices}
\begin{solution}
  Unions, intersections, tuples, and primitives are commonly aliased with \lstinline[style=tsinline]|type|. Interfaces cannot express a bare union of primitives as a single interface body.
\end{solution}

\question[4]
A \textbf{function type alias} can be written as:
\begin{choices}
  \choice \lstinline[style=tsinline]|type Compare = function(a: number, b: number): number;|
  \CorrectChoice \lstinline[style=tsinline]|type Compare = (a: number, b: number) => number;|
  \choice \lstinline[style=tsinline]|interface Compare() => number;|
  \choice \lstinline[style=tsinline]|type Compare = number => number => number;| without parentheses
\end{choices}
\begin{solution}
\begin{lstlisting}[style=tsexam]
type Compare = (a: number, b: number) => number;
type Handler = (event: string) => void;
let sort: Compare = (x, y) => x - y;
\end{lstlisting}
\end{solution}

\question[4]
Type aliases are especially useful for:
\begin{choices}
  \choice Replacing all JavaScript functions with SQL
  \CorrectChoice Naming complex types (unions, tuples, function shapes) and reusing them across a codebase
  \choice Creating runtime singleton objects automatically
  \choice Disabling \texttt{strict} in \texttt{tsconfig.json}
\end{choices}
\begin{solution}
  Example: \lstinline[style=tsinline]|type UserId = string;| documents intent better than repeating \lstinline[style=tsinline]|string| everywhere. Change the alias once if the model evolves.
\end{solution}

\question[4]
Compared with \textbf{interfaces}, type aliases can:
\begin{choices}
  \choice Merge across files by default when two modules declare the same name
  \CorrectChoice Represent unions, intersections, tuples, and mapped/conditional types that interfaces cannot express directly
  \choice Only alias primitive \lstinline[style=tsinline]|string|
  \choice Emit extra runtime classes
\end{choices}
\begin{solution}
  Interfaces support \lstinline[style=tsinline]|extends| and declaration merging for object shapes. Use \lstinline[style=tsinline]|type| when the right-hand side is not a plain extendable object shape (see also Part~7).
\end{solution}

\question[4]
Given \lstinline[style=tsinline]|type Point = { x: number; y: number };|, assigning \lstinline[style=tsinline]|{ x: 1, y: 2, z: 3 }| to \lstinline[style=tsinline]|Point|:
\begin{choices}
  \choice Is always allowed (extra properties always OK on variables)
  \CorrectChoice May error on the object literal directly (excess property checking) unless a wider variable type is used
  \choice Requires \lstinline[style=tsinline]|enum|
  \choice Converts \lstinline[style=tsinline]|z| to \lstinline[style=tsinline]|void|
\end{choices}
\begin{solution}
  Fresh object literals are checked strictly; assign to a variable typed \lstinline[style=tsinline]|Point| first or use a type assertion only when you are certain (Part~21).
\end{solution}

\uplevel{\par\textbf{Part 20: Interfaces (object contracts)}\par}

\question[4]
An \textbf{interface} describes:
\begin{choices}
  \choice A compiled \texttt{.js} file on disk
  \CorrectChoice The shape of an object: required/optional properties and method signatures
  \choice Only CSS classes in Angular components
  \choice npm package versions
\end{choices}
\begin{solution}
  Interfaces exist only at compile time. They guide autocomplete and catch missing or mistyped properties when you create or consume objects.
\end{solution}

\question[4]
One interface extending another uses:
\begin{choices}
  \choice \lstinline[style=tsinline]|interface B implements A|
  \CorrectChoice \lstinline[style=tsinline]|interface B extends A { ... }|
  \choice \lstinline[style=tsinline]|interface B : A|
  \choice \lstinline[style=tsinline]|extends interface B from A|
\end{choices}
\begin{solution}
\begin{lstlisting}[style=tsexam]
interface Person { name: string; }
interface Employee extends Person {
  employeeId: number;
}
\end{lstlisting}
  A class uses \lstinline[style=tsinline]|implements| to satisfy interfaces; interfaces use \lstinline[style=tsinline]|extends| to inherit shapes.
\end{solution}

\question[4]
A class may \lstinline[style=tsinline]|implements| multiple interfaces:
\begin{choices}
  \choice False---only one interface per class
  \CorrectChoice True---\lstinline[style=tsinline]|class C implements A, B { }| must satisfy every listed contract
  \choice True only if both interfaces are empty
  \choice False---use \lstinline[style=tsinline]|extends| instead of \lstinline[style=tsinline]|implements|
\end{choices}
\begin{solution}
  \lstinline[style=tsinline]|implements| checks structural compatibility at compile time. It does not inherit implementation from interfaces---you still write method bodies in the class.
\end{solution}

\question[4]
An interface with a \textbf{call signature}:
\begin{lstlisting}[style=tsexam]
interface Formatter {
  (value: string): string;
}
\end{lstlisting}
describes:
\begin{choices}
  \choice An array of strings only
  \CorrectChoice An object type that is also callable like a function
  \choice A CSS module
  \choice A route guard
\end{choices}
\begin{solution}
  Callable object types model functions with extra properties (\lstinline[style=tsinline]|fn.props = ...|). Function type aliases are simpler when you only need a plain function shape.
\end{solution}

\question[4]
\textbf{Declaration merging} allows:
\begin{choices}
  \choice Two \lstinline[style=tsinline]|type| aliases with the same name to merge automatically
  \CorrectChoice Multiple \lstinline[style=tsinline]|interface| declarations with the same name to combine members (common in ambient \texttt{.d.ts} libraries)
  \choice Merging \lstinline[style=tsinline]|class| bodies across files
  \choice Replacing \texttt{tsconfig.json} with \texttt{package.json}
\end{choices}
\begin{solution}
  Example: extending global \lstinline[style=tsinline]|Window| in a \texttt{.d.ts} file. Type aliases do not merge---duplicate \lstinline[style=tsinline]|type| names are errors.
\end{solution}

\question[4]
An interface property \lstinline[style=tsinline]|readonly name: string| means:
\begin{choices}
  \choice The property may be reassigned after object creation
  \CorrectChoice Assignments to \lstinline[style=tsinline]|name| after creation are disallowed (compile-time)
  \choice The property is optional
  \choice The interface cannot be implemented by a class
\end{choices}
\begin{solution}
  \lstinline[style=tsinline]|ReadonlyArray<T>| and \lstinline[style=tsinline]|readonly| fields document immutability intent. Runtime freezing still requires \lstinline[style=tsinline]|Object.freeze| if you need enforcement there.
\end{solution}

\uplevel{\par\textbf{Part 21: Type casting and assertions}\par}

\question[4]
The TypeScript \textbf{assertion} syntax \lstinline[style=tsinline]|value as string|:
\begin{choices}
  \choice Converts the runtime value to a string automatically
  \CorrectChoice Tells the compiler to treat \lstinline[style=tsinline]|value| as \lstinline[style=tsinline]|string| without changing runtime behavior
  \choice Deletes \lstinline[style=tsinline]|value| from memory
  \choice Enables \texttt{strict} mode
\end{choices}
\begin{solution}
  Assertions are compile-time only. They do not validate data---incorrect assertions can hide bugs that surface at runtime.
\end{solution}

\question[4]
After narrowing fails, casting with \lstinline[style=tsinline]|as any| and then using the value:
\begin{choices}
  \choice Is the safest pattern for all external data
  \CorrectChoice Bypasses type checking and should be rare---prefer validation or \lstinline[style=tsinline]|unknown| plus narrowing
  \choice Emits stronger runtime checks automatically
  \choice Replaces \lstinline[style=tsinline]|interface| declarations
\end{choices}
\begin{solution}
  Double casts (\lstinline[style=tsinline]|x as any as Target|) are a code smell. Use type guards, schema validation (zod/io-ts), or proper generics when integrating untyped APIs.
\end{solution}

\question[4]
The \textbf{non-null assertion operator} \lstinline[style=tsinline]|user!.name| means:
\begin{choices}
  \choice \lstinline[style=tsinline]|user| is definitely \lstinline[style=tsinline]|null| at runtime
  \CorrectChoice You assert \lstinline[style=tsinline]|user| is not \lstinline[style=tsinline]|null| or \lstinline[style=tsinline]|undefined| (compiler only; runtime may still throw)
  \choice \lstinline[style=tsinline]|name| is optional always
  \choice The property is \lstinline[style=tsinline]|readonly|
\end{choices}
\begin{solution}
  Use sparingly after checks you know are true (for example right after \lstinline[style=tsinline]|if (!user) return;|). Prefer optional chaining \lstinline[style=tsinline]|user?.name| when unsure.
\end{solution}

\question[4]
\texttt{const} assertion \lstinline[style=tsinline]|as const| on an object:
\begin{choices}
  \choice Widens all properties to \lstinline[style=tsinline]|any|
  \CorrectChoice Makes properties readonly literal types (for example \lstinline[style=tsinline]|'left'| instead of \lstinline[style=tsinline]|string|)
  \choice Converts the object to an \lstinline[style=tsinline]|enum|
  \choice Disables \lstinline[style=tsinline]|strictNullChecks|
\end{choices}
\begin{solution}
\begin{lstlisting}[style=tsexam]
const config = { mode: 'dark', level: 1 } as const;
// config.mode is type 'dark', not string
\end{lstlisting}
  Useful for constant lookup tables and discriminant maps.
\end{solution}

\question[4]
Casting is \textbf{preferable} to narrowing when:
\begin{choices}
  \choice You have not verified external data and want safety
  \CorrectChoice You already know more than the checker (for example DOM APIs: \lstinline[style=tsinline]|document.getElementById('x') as HTMLInputElement|) and will validate separately if needed
  \choice You want the compiler to prove exhaustiveness
  \choice You are assigning a literal union without checks
\end{choices}
\begin{solution}
  Prefer narrowing (\lstinline[style=tsinline]|if (el instanceof HTMLInputElement)|) when possible. Assertions document trust; they do not replace runtime validation for user input.
\end{solution}

\question[4]
Legacy angle-bracket syntax \lstinline[style=tsinline]|<string>value|:
\begin{choices}
  \choice Is required in all \texttt{.tsx} files
  \CorrectChoice Is equivalent to \lstinline[style=tsinline]|value as string| but cannot be used in \texttt{.tsx} because \texttt{<>} conflicts with JSX
  \choice Runs only in \texttt{strict} mode
  \choice Emits a runtime cast in JavaScript
\end{choices}
\begin{solution}
  In \texttt{.tsx}, use \lstinline[style=tsinline]|as| form. Both assertion styles disappear in emitted JavaScript.
\end{solution}

\newpage

\uplevel{\par\textbf{Part 22: Functions}\par}

\question[4]
A function parameter type annotation is written as:
\begin{choices}
  \choice \lstinline[style=tsinline]|function add(a: number): string|
  \CorrectChoice \lstinline[style=tsinline]|function add(a: number, b: number): number|
  \choice \lstinline[style=tsinline]|function add: (a, b) => number|
  \choice \lstinline[style=tsinline]|function add[a number]|
\end{choices}
\begin{solution}
  Syntax: \lstinline[style=tsinline]|function name(param: Type, ...): ReturnType|. Return type inference works for simple cases but annotate public APIs for clarity.
\end{solution}

\question[4]
An \textbf{optional parameter} in TypeScript is marked with:
\begin{choices}
  \choice \lstinline[style=tsinline]|@Optional|
  \CorrectChoice \lstinline[style=tsinline]|?| after the name (\lstinline[style=tsinline]|name?: string|)
  \choice \lstinline[style=tsinline]|...| before the name
  \choice \lstinline[style=tsinline]|readonly|
\end{choices}
\begin{solution}
  Callers may omit the argument; inside the function it is \lstinline[style=tsinline]|string | undefined|. Default parameters (\lstinline[style=tsinline]|size = 10|) supply a value when the argument is \lstinline[style=tsinline]|undefined|.
\end{solution}

\question[4]
A \textbf{rest parameter} collecting extra arguments uses:
\begin{choices}
  \choice \lstinline[style=tsinline]|function sum(a, b, c)|
  \CorrectChoice \lstinline[style=tsinline]|function sum(...nums: number[])| or tuple rest types
  \choice \lstinline[style=tsinline]|function sum[number]|
  \choice \lstinline[style=tsinline]|function sum?: number[]|
\end{choices}
\begin{solution}
\begin{lstlisting}[style=tsexam]
function sum(...nums: number[]): number {
  return nums.reduce((a, b) => a + b, 0);
}
\end{lstlisting}
  Rest must be last in the parameter list. Typed as an array or tuple for fixed positions.
\end{solution}

\question[4]
Which arrow function signature is correctly typed?
\begin{choices}
  \choice \lstinline[style=tsinline]|const f: number = (x) => x * 2;|
  \CorrectChoice \lstinline[style=tsinline]|const f = (x: number): number => x * 2;| or \lstinline[style=tsinline]|const f: (x: number) => number = (x) => x * 2;|
  \choice \lstinline[style=tsinline]|const f => (x: number) => number;|
  \choice \lstinline[style=tsinline]|arrow function f(x: number): number|
\end{choices}
\begin{solution}
  Arrow functions inherit \lstinline[style=tsinline]|this| from the enclosing scope (unlike classic \lstinline[style=tsinline]|function|), which matters in callbacks and Angular component code.
\end{solution}

\question[4]
\textbf{Function overload signatures} allow:
\begin{choices}
  \choice Only one call signature per function name ever
  \CorrectChoice Multiple call signatures for one implementation (for example accept \lstinline[style=tsinline]|string| or \lstinline[style=tsinline]|string[]| with different return types)
  \choice Replacing \lstinline[style=tsinline]|return| statements
  \choice Runtime polymorphism like Java virtual methods automatically
\end{choices}
\begin{solution}
\begin{lstlisting}[style=tsexam]
function format(input: string): string;
function format(input: string[]): string[];
function format(input: string | string[]): string | string[] {
  return Array.isArray(input) ? input.map(s => s.trim()) : input.trim();
}
\end{lstlisting}
  Overloads are a compile-time API surface; one implementation body handles all cases.
\end{solution}

\question[4]
A function with no \texttt{return} statement and return type \lstinline[style=tsinline]|void|:
\begin{choices}
  \choice Must return \lstinline[style=tsinline]|null| always
  \CorrectChoice May return \lstinline[style=tsinline]|undefined| implicitly; callers should ignore the result
  \choice Cannot be arrow functions
  \choice Must throw \lstinline[style=tsinline]|never|
\end{choices}
\begin{solution}
  Annotate side-effect functions (\lstinline[style=tsinline]|log|, \lstinline[style=tsinline]|saveToDisk|) with \lstinline[style=tsinline]|void|. Do not use \lstinline[style=tsinline]|void| as a generic ``ignore anything'' parameter type---use \lstinline[style=tsinline]|unknown| instead.
\end{solution}

\question[4]
A default parameter \lstinline[style=tsinline]|function greet(name = 'Guest')| means:
\begin{choices}
  \choice \lstinline[style=tsinline]|name| must always be passed explicitly
  \CorrectChoice Callers may omit \lstinline[style=tsinline]|name|; it defaults to \lstinline[style=tsinline]|'Guest'| and is inferred as \lstinline[style=tsinline]|string|
  \choice \lstinline[style=tsinline]|name| is optional and always \lstinline[style=tsinline]|undefined| when omitted
  \choice The function cannot return a value
\end{choices}
\begin{solution}
  Default parameters are not the same as optional parameters (\lstinline[style=tsinline]|name?: string|), but both allow omission with different typing rules inside the body.
\end{solution}

\newpage

\uplevel{\par\textbf{Part 23: Classes (inheritance and abstraction)}\par}

\question[4]
A TypeScript \textbf{class} combines:
\begin{choices}
  \choice Only HTML templates and CSS
  \CorrectChoice Fields (state), a constructor, and methods; it may \lstinline[style=tsinline]|implement| interfaces and \lstinline[style=tsinline]|extend| other classes
  \choice Only global variables on \lstinline[style=tsinline]|window|
  \choice SQL table definitions only
\end{choices}
\begin{solution}
\begin{lstlisting}[style=tsexam]
class Counter {
  count = 0;
  increment() { this.count += 1; }
}
\end{lstlisting}
  Classes compile to JavaScript prototype or class syntax depending on \texttt{target}. They support the same OOP patterns as other class-based languages.
\end{solution}

\question[4]
\textbf{Inheritance} in TypeScript uses the keyword:
\begin{choices}
  \choice \lstinline[style=tsinline]|implements| on the parent class
  \CorrectChoice \lstinline[style=tsinline]|extends| (\lstinline[style=tsinline]|class Dog extends Animal { }|)
  \choice \lstinline[style=tsinline]|inherits|
  \choice \lstinline[style=tsinline]|superclass|
\end{choices}
\begin{solution}
  \lstinline[style=tsinline]|extends| copies the prototype chain and brings inherited methods. \lstinline[style=tsinline]|implements| satisfies a type contract without inheriting implementation.
\end{solution}

\question[4]
In a subclass constructor, \lstinline[style=tsinline]|super(...)| must be called:
\begin{choices}
  \choice Never---subclasses have no constructors
  \CorrectChoice Before using \lstinline[style=tsinline]|this| when the base class defines a constructor (to run parent initialization)
  \choice Only in abstract classes
  \choice After every \lstinline[style=tsinline]|return| statement
\end{choices}
\begin{solution}
\begin{lstlisting}[style=tsexam]
class Animal { constructor(public name: string) {} }
class Dog extends Animal {
  constructor(name: string, public breed: string) {
    super(name);
  }
}
\end{lstlisting}
\end{solution}

\question[4]
An \textbf{abstract class}:
\begin{choices}
  \choice Can be instantiated directly with \lstinline[style=tsinline]|new| like any class
  \CorrectChoice Cannot be instantiated directly; it may declare abstract members subclasses must implement
  \choice Must contain only static methods
  \choice Replaces interfaces entirely
\end{choices}
\begin{solution}
\begin{lstlisting}[style=tsexam]
abstract class Shape {
  abstract area(): number;
}
class Circle extends Shape {
  constructor(public r: number) { super(); }
  area() { return Math.PI * this.r ** 2; }
}
\end{lstlisting}
  Use abstraction to share code in the base class while forcing concrete subclasses to fill in behavior.
\end{solution}

\question[4]
The \lstinline[style=tsinline]|override| keyword (with \texttt{noImplicitOverride}) helps:
\begin{choices}
  \choice Delete methods from the parent class at runtime
  \CorrectChoice Ensure a method actually overrides a base member (typo-safe when parents rename methods)
  \choice Mark fields as \lstinline[style=tsinline]|readonly|
  \choice Enable decorators automatically
\end{choices}
\begin{solution}
  \lstinline[style=tsinline]|override greet() { }| documents intent and catches mistakes such as misspelling the base method name when the base class changes.
\end{solution}

\question[4]
\textbf{Static} members on a class:
\begin{choices}
  \choice Belong to each instance and use \lstinline[style=tsinline]|this| only
  \CorrectChoice Belong to the constructor function (shared by all instances), accessed as \lstinline[style=tsinline]|ClassName.member|
  \choice Cannot be methods
  \choice Replace \lstinline[style=tsinline]|interface| merging
\end{choices}
\begin{solution}
\begin{lstlisting}[style=tsexam]
class MathUtil {
  static PI = 3.14159;
  static clamp(n: number, min: number, max: number) {
    return Math.min(max, Math.max(min, n));
  }
}
\end{lstlisting}
  Static members are not on instances---useful for factories and constants, but avoid overusing static state in apps.
\end{solution}

\question[4]
\textbf{Abstraction} in OOP with TypeScript typically means:
\begin{choices}
  \choice Hiding all types with \lstinline[style=tsinline]|any|
  \CorrectChoice Exposing essential behavior through classes/interfaces while hiding implementation details (abstract classes, interfaces, private fields)
  \choice Removing constructors from every class
  \choice Compiling without \texttt{tsconfig.json}
\end{choices}
\begin{solution}
  Callers depend on \lstinline[style=tsinline]|Shape.area()| or an \lstinline[style=tsinline]|interface| contract, not internal fields. Angular services and components use the same idea at framework scale.
\end{solution}

\uplevel{\par\textbf{Part 24: Access modifiers}\par}

\question[4]
In TypeScript classes, \lstinline[style=tsinline]|public| members (the default):
\begin{choices}
  \choice Are visible only inside the declaring file
  \CorrectChoice Are accessible from anywhere code may reference the instance (unless type system limits the reference)
  \choice Cannot include methods
  \choice Are the same as \lstinline[style=tsinline]|protected|
\end{choices}
\begin{solution}
  \lstinline[style=tsinline]|public name: string| in the constructor parameter list creates a public field automatically (parameter property).
\end{solution}

\question[4]
A \lstinline[style=tsinline]|protected| field is accessible:
\begin{choices}
  \choice Only inside the same source file, never in subclasses
  \CorrectChoice Inside the declaring class and its subclasses, but not from unrelated outside code
  \choice Only from static methods
  \choice Only through \lstinline[style=tsinline]|JSON.parse|
\end{choices}
\begin{solution}
  Use \lstinline[style=tsinline]|protected| for extension points in base classes. Subclasses can read/write protected members; consumers of the subclass should use public API instead.
\end{solution}

\question[4]
\lstinline[style=tsinline]|private| on a class field means:
\begin{choices}
  \choice The field is encrypted at runtime
  \CorrectChoice Only the declaring class body may access it (compile-time check; erased in JS emit)
  \choice Subclasses always inherit access to the field
  \choice The field must be \lstinline[style=tsinline]|static|
\end{choices}
\begin{solution}
  TypeScript \lstinline[style=tsinline]|private| is not true runtime privacy. ECMAScript \lstinline[style=tsinline]|#count| fields are private at runtime. Many codebases use TS \lstinline[style=tsinline]|private| for clarity inside class files.
\end{solution}

\question[4]
\textbf{Parameter properties} in a constructor:
\begin{lstlisting}[style=tsexam]
constructor(private id: number, public name: string) {}
\end{lstlisting}
\begin{choices}
  \choice Declare parameters only---no fields are created
  \CorrectChoice Declare and initialize class fields in one step from constructor parameters
  \choice Require \lstinline[style=tsinline]|abstract| on the class
  \choice Work only in interfaces
\end{choices}
\begin{solution}
  Equivalent to declaring fields and assigning in the constructor body. Common in Angular services and components to reduce boilerplate.
\end{solution}

\question[4]
A \lstinline[style=tsinline]|readonly| constructor parameter \lstinline[style=tsinline]|constructor(readonly id: number)|:
\begin{choices}
  \choice Allows reassignment to \lstinline[style=tsinline]|id| anywhere in the class after construction
  \CorrectChoice Creates a field assignable only during initialization (typically in the constructor)
  \choice Makes the class \lstinline[style=tsinline]|abstract|
  \choice Exports the field to \texttt{package.json}
\end{choices}
\begin{solution}
  Pair with \lstinline[style=tsinline]|private|/\lstinline[style=tsinline]|public|/\lstinline[style=tsinline]|protected| as needed: \lstinline[style=tsinline]|constructor(private readonly store: Store)|.
\end{solution}

\question[4]
From outside the class, accessing a \lstinline[style=tsinline]|private| method:
\begin{choices}
  \choice Is always allowed in JavaScript and TypeScript
  \CorrectChoice Causes a TypeScript compile error (even if JavaScript might still allow it at runtime depending on emit)
  \choice Requires \lstinline[style=tsinline]|as protected|
  \choice Requires the \lstinline[style=tsinline]|override| keyword on the caller
\end{choices}
\begin{solution}
  Encapsulation keeps invariants inside the class. Expose behavior through \lstinline[style=tsinline]|public| methods rather than reaching into \lstinline[style=tsinline]|private| state from collaborators.
\end{solution}

\uplevel{\par\textbf{Part 25: Decorators}\par}

\question[4]
In TypeScript, a \textbf{decorator} is:
\begin{choices}
  \choice A CSS class applied with \lstinline[style=htmlexam]|[ngClass]|
  \CorrectChoice A special declaration (often a function) attached to a class, method, property, or parameter that can modify or annotate that declaration
  \choice A replacement for \lstinline[style=tsinline]|interface|
  \choice A SQL view on a database table
\end{choices}
\begin{solution}
  Decorators are enabled with \texttt{experimentalDecorators} (and often \texttt{emitDecoratorMetadata}) in \texttt{tsconfig.json}. Angular relies on this model for \lstinline[style=tsinline]|@Component|, \lstinline[style=tsinline]|@Injectable|, and others.
\end{solution}

\question[4]
A \textbf{class decorator} such as \lstinline[style=tsinline]|@Component({...})|:
\begin{choices}
  \choice Runs only in the browser after click events
  \CorrectChoice Receives the class constructor and can register or modify the class (Angular registers metadata for DI and compilation)
  \choice Replaces the need for a \texttt{templateUrl}
  \choice Must be placed on every local variable
\end{choices}
\begin{solution}
  Class decorators are functions evaluated at class definition time. Frameworks read metadata to wire templates, selectors, and providers.
\end{solution}

\question[4]
A \textbf{method decorator} is applied to:
\begin{choices}
  \choice A function parameter in the parameter list
  \CorrectChoice A method inside a class (for example logging, validation, or framework hooks)
  \choice An \lstinline[style=tsinline]|import| statement
  \choice A \texttt{tsconfig.json} file
\end{choices}
\begin{solution}
  Method decorators receive information about the method (name, descriptor). Use cases include cross-cutting concerns; Angular uses decorators heavily at class level rather than ad-hoc method decorators in app code.
\end{solution}

\question[4]
A \textbf{property decorator} such as \lstinline[style=tsinline]|@Input()| or \lstinline[style=tsinline]|@Output()|:
\begin{choices}
  \choice Declares a route path in the router
  \CorrectChoice Marks a class field for framework processing (for example Angular input/output binding metadata)
  \choice Makes the property \lstinline[style=tsinline]|static| automatically
  \choice Replaces TypeScript type annotations
\end{choices}
\begin{solution}
  Property decorators run against the prototype property (or field). \lstinline[style=tsinline]|@Input() title: string| tells Angular the parent may bind \lstinline[style=htmlexam]|[title]|; \lstinline[style=tsinline]|@Output() saved| pairs with an \lstinline[style=tsinline]|EventEmitter|.
\end{solution}

\question[4]
A \textbf{parameter decorator} such as \lstinline[style=tsinline]|@Inject(TOKEN)| is used to:
\begin{choices}
  \choice Style HTML elements with CSS
  \CorrectChoice Attach metadata to a constructor/method parameter (Angular DI uses this to know what to inject)
  \choice Convert \lstinline[style=tsinline]|string| to \lstinline[style=tsinline]|number| at runtime
  \choice Enable \lstinline[style=tsinline]|strict| mode
\end{choices}
\begin{solution}
  Parameter decorators do not change the parameter value by themselves; they record metadata read by the injector. Many apps use type-based injection without explicit \lstinline[style=tsinline]|@Inject| when types are available.
\end{solution}

\question[4]
Which pair belongs to Angular \textbf{dependency injection} and decorators?
\begin{choices}
  \choice \lstinline[style=tsinline]|@Pipe| on a variable and \lstinline[style=tsinline]|*ngIf|
  \CorrectChoice \lstinline[style=tsinline]|@Injectable()| on a service class and constructor injection of that service
  \choice \lstinline[style=tsinline]|@NgModule| on a local \lstinline[style=tsinline]|let| variable
  \choice \lstinline[style=tsinline]|@Component| on \texttt{package.json}
\end{choices}
\begin{solution}
\begin{lstlisting}[style=tsexam]
@Injectable({ providedIn: 'root' })
export class DataService { }

@Component({ selector: 'app-root', template: `...` })
export class AppComponent {
  constructor(private data: DataService) {}
}
\end{lstlisting}
  Decorators supply metadata; the injector uses it at runtime in the compiled app.
\end{solution}

\question[4]
To use legacy/experimental decorators in a project, \texttt{tsconfig.json} typically sets:
\begin{choices}
  \choice \texttt{"target": "ES3"} only
  \CorrectChoice \texttt{"experimentalDecorators": true} (and often \texttt{"emitDecoratorMetadata": true} for Angular)
  \choice \texttt{"strict": false} only
  \choice \texttt{"module": "CSS"}
\end{choices}
\begin{solution}
  TypeScript 5+ introduced standard decorators with different configuration; Angular toolchains still commonly document the experimental flags. Match your CLI version's recommended \texttt{tsconfig} settings.
\end{solution}

\newpage

\uplevel{\par\textbf{Part 26: Utility types}\par}

\question[4]
TypeScript \textbf{utility types} are:
\begin{choices}
  \choice Runtime helper classes shipped in every browser
  \CorrectChoice Built-in generic type transformations (such as \lstinline[style=tsinline]|Partial<T>|) defined in the standard library
  \choice CSS modules for Angular components
  \choice npm scripts in \texttt{package.json}
\end{choices}
\begin{solution}
  Utilities compose existing types without duplicating field lists. They are compile-time only and erase when emitting JavaScript.
\end{solution}

\question[4]
\lstinline[style=tsinline]|Partial<User>| when \lstinline[style=tsinline]|User| has required fields means:
\begin{choices}
  \choice Every property becomes required and non-nullable
  \CorrectChoice Every property of \lstinline[style=tsinline]|User| becomes optional
  \choice Only string fields remain
  \choice The type becomes \lstinline[style=tsinline]|never|
\end{choices}
\begin{solution}
  Useful for update DTOs (\lstinline[style=tsinline]|patchUser(id, partial: Partial<User>)|) where callers send only changed fields.
\end{solution}

\question[4]
\lstinline[style=tsinline]|Required<Config>| compared with \lstinline[style=tsinline]|Config| when some fields were optional:
\begin{choices}
  \choice Makes all properties optional
  \CorrectChoice Makes every property required (removes optionality from \lstinline[style=tsinline]|?| fields)
  \choice Deletes properties at runtime
  \choice Converts the type to a union only
\end{choices}
\begin{solution}
  Pair with \lstinline[style=tsinline]|Partial| when building defaults: accept \lstinline[style=tsinline]|Partial<Config>|, merge defaults, then treat as \lstinline[style=tsinline]|Required<Config>| internally.
\end{solution}

\question[4]
\lstinline[style=tsinline]|Pick<User, 'id' | 'name'>| produces a type with:
\begin{choices}
  \choice All keys of \lstinline[style=tsinline]|User| except \lstinline[style=tsinline]|id| and \lstinline[style=tsinline]|name|
  \CorrectChoice Only the listed properties \lstinline[style=tsinline]|id| and \lstinline[style=tsinline]|name| (with their original types)
  \choice No properties---empty object type always
  \choice Only methods, never fields
\end{choices}
\begin{solution}
  \lstinline[style=tsinline]|Omit<User, 'password'>| is the opposite---everything but the omitted keys. Both are built from \lstinline[style=tsinline]|keyof| (Part~28).
\end{solution}

\question[4]
\lstinline[style=tsinline]|Readonly<T>| on an interface:
\begin{choices}
  \choice Allows assigning new values to every property freely
  \CorrectChoice Makes all properties readonly (cannot assign to properties after creation)
  \choice Converts numbers to strings
  \choice Emits runtime freeze code automatically always
\end{choices}
\begin{solution}
  Similar to marking each field \lstinline[style=tsinline]|readonly|. For arrays use \lstinline[style=tsinline]|ReadonlyArray<T>| or \lstinline[style=tsinline]|readonly T[]|.
\end{solution}

\question[4]
\lstinline[style=tsinline]|ReturnType<typeof fn>| extracts:
\begin{choices}
  \choice The parameter types of \lstinline[style=tsinline]|fn| as a tuple
  \CorrectChoice The return type of function \lstinline[style=tsinline]|fn|
  \choice The \lstinline[style=tsinline]|this| type of a class
  \choice The npm name of a package
\end{choices}
\begin{solution}
  Related utilities: \lstinline[style=tsinline]|Parameters<typeof fn>| for argument tuple types, \lstinline[style=tsinline]|Awaited<PromiseType>| for promise results. Handy when wrapping library functions without repeating types.
\end{solution}

\question[4]
\lstinline[style=tsinline]|Exclude<'a' | 'b' | 'c', 'a'>| evaluates to:
\begin{choices}
  \choice \lstinline[style=tsinline]|'a'| only
  \CorrectChoice \lstinline[style=tsinline]|'b' | 'c'|
  \choice \lstinline[style=tsinline]|never| always
  \choice \lstinline[style=tsinline]|string| always
\end{choices}
\begin{solution}
  \lstinline[style=tsinline]|Extract<T, U>| keeps union members assignable to \lstinline[style=tsinline]|U|. \lstinline[style=tsinline]|NonNullable<T>| removes \lstinline[style=tsinline]|null| and \lstinline[style=tsinline]|undefined| from \lstinline[style=tsinline]|T|.
\end{solution}

\uplevel{\par\textbf{Part 27: Array generics}\par}

\question[4]
A generic array type for strings can be written as:
\begin{choices}
  \choice \lstinline[style=tsinline]|Array strings|
  \CorrectChoice \lstinline[style=tsinline]|Array<string>| or \lstinline[style=tsinline]|string[]|
  \choice \lstinline[style=tsinline]|string<>|
  \choice \lstinline[style=tsinline]|[string]| only (always a one-tuple)
\end{choices}
\begin{solution}
  \lstinline[style=tsinline]|T[]| is syntax sugar for \lstinline[style=tsinline]|Array<T>|. Both allow any length (unlike fixed tuples in Part~15).
\end{solution}

\question[4]
A generic function \lstinline[style=tsinline]|function first<T>(items: T[]): T | undefined|:
\begin{choices}
  \choice Requires \lstinline[style=tsinline]|T| to be \lstinline[style=tsinline]|string| only
  \CorrectChoice Returns the same element type as the array elements (or \lstinline[style=tsinline]|undefined| if empty)
  \choice Returns \lstinline[style=tsinline]|any| always
  \choice Cannot accept empty arrays at compile time
\end{choices}
\begin{solution}
  One type parameter keeps input and output related: \lstinline[style=tsinline]|first([1,2,3])| is \lstinline[style=tsinline]|number | undefined|; \lstinline[style=tsinline]|first(['a'])| is \lstinline[style=tsinline]|string | undefined|.
\end{solution}

\question[4]
Calling \lstinline[style=tsinline]|numbers.map(n => n * 2)| when \lstinline[style=tsinline]|numbers: number[]| infers:
\begin{choices}
  \choice \lstinline[style=tsinline]|map| returns \lstinline[style=tsinline]|any[]| always
  \CorrectChoice The callback parameter \lstinline[style=tsinline]|n| is \lstinline[style=tsinline]|number| and the result is \lstinline[style=tsinline]|number[]|
  \choice \lstinline[style=tsinline]|n| is \lstinline[style=tsinline]|string|
  \choice The array becomes a tuple of length 1 only
\end{choices}
\begin{solution}
  \lstinline[style=tsinline]|Array<T>| methods are generic: \lstinline[style=tsinline]|filter|, \lstinline[style=tsinline]|map|, \lstinline[style=tsinline]|reduce| preserve or transform element types when callbacks are typed.
\end{solution}

\question[4]
\lstinline[style=tsinline]|ReadonlyArray<number>| differs from \lstinline[style=tsinline]|number[]| because:
\begin{choices}
  \choice It allows \lstinline[style=tsinline]|push| and \lstinline[style=tsinline]|pop|
  \CorrectChoice Mutating methods such as \lstinline[style=tsinline]|push| are unavailable on the readonly type
  \choice It cannot be iterated
  \choice It stores only one number
\end{choices}
\begin{solution}
  Use readonly arrays when exposing internal data you do not want consumers to mutate. \lstinline[style=tsinline]|readonly number[]| is a similar shorthand.
\end{solution}

\question[4]
A \textbf{generic constraint} on an array parameter \lstinline[style=tsinline]|T extends { id: number }| ensures:
\begin{choices}
  \choice \lstinline[style=tsinline]|T| must be a primitive
  \CorrectChoice Elements have at least an \lstinline[style=tsinline]|id: number| field so \lstinline[style=tsinline]|item.id| is legal inside the function
  \choice The array length is exactly 1
  \choice \lstinline[style=tsinline]|map| is disabled
\end{choices}
\begin{solution}
\begin{lstlisting}[style=tsexam]
function findById<T extends { id: number }>(items: T[], id: number): T | undefined {
  return items.find(item => item.id === id);
}
\end{lstlisting}
\end{solution}

\question[4]
\lstinline[style=tsinline]|Array.isArray(value)| with a generic \lstinline[style=tsinline]|unknown| value:
\begin{choices}
  \choice Narrows \lstinline[style=tsinline]|value| to \lstinline[style=tsinline]|string[]| automatically
  \CorrectChoice Narrows to \lstinline[style=tsinline]|any[]| at compile time (element type still unknown---further checks needed)
  \choice Proves every element is a number
  \choice Replaces the need for tuples
\end{choices}
\begin{solution}
  After \lstinline[style=tsinline]|Array.isArray(value)|, you know it is an array but not \lstinline[style=tsinline]|T[]| until you validate elements or cast carefully. Prefer typed APIs over repeated casting.
\end{solution}

\uplevel{\par\textbf{Part 28: \texttt{keyof} and indexed access}\par}

\question[4]
For an object type \lstinline[style=tsinline]|User|, the type \lstinline[style=tsinline]|keyof User| is:
\begin{choices}
  \choice The runtime list of values in the object
  \CorrectChoice A union of that type's property names as string literal types (for example \lstinline[style=tsinline]|'id' | 'name'|)
  \choice Always \lstinline[style=tsinline]|string|
  \choice A numeric array length
\end{choices}
\begin{solution}
\begin{lstlisting}[style=tsexam]
interface User { id: number; name: string; }
type UserKeys = keyof User;  // 'id' | 'name'
\end{lstlisting}
  See also Part~9 \lstinline[style=tsinline]|getProperty| example.
\end{solution}

\question[4]
Given \lstinline[style=tsinline]|type K = keyof { a: 1; b: 2 }|, \lstinline[style=tsinline]|K| is equivalent to:
\begin{choices}
  \choice \lstinline[style=tsinline]|number|
  \CorrectChoice \lstinline[style=tsinline]|'a' | 'b'|
  \choice \lstinline[style=tsinline]|1 | 2|
  \choice \lstinline[style=tsinline]|object|
\end{choices}
\begin{solution}
  \lstinline[style=tsinline]|keyof| operates on \textbf{types}, not runtime objects. For values, use \lstinline[style=tsinline]|Object.keys| (returns \lstinline[style=tsinline]|string[]| unless narrowed).
\end{solution}

\question[4]
An \textbf{indexed access type} \lstinline[style=tsinline]|User['name']| resolves to:
\begin{choices}
  \choice The literal type \lstinline[style=tsinline]|'name'|
  \CorrectChoice The type of the \lstinline[style=tsinline]|name| property (\lstinline[style=tsinline]|string| in a typical \lstinline[style=tsinline]|User|)
  \choice \lstinline[style=tsinline]|never|
  \choice A union of all keys
\end{choices}
\begin{solution}
  \lstinline[style=tsinline]|User[keyof User]| is a union of all property value types. Useful for generic getters and mapped types.
\end{solution}

\question[4]
\lstinline[style=tsinline]|Pick<T, K extends keyof T>| requires \lstinline[style=tsinline]|K| to be:
\begin{choices}
  \choice Any arbitrary \lstinline[style=tsinline]|string| with no check
  \CorrectChoice A subset of keys of \lstinline[style=tsinline]|T| (enforced by \lstinline[style=tsinline]|extends keyof T|)
  \choice Only numeric keys
  \choice A runtime array passed to \lstinline[style=tsinline]|JSON.parse|
\end{choices}
\begin{solution}
  \lstinline[style=tsinline]|Pick<User, 'id'>| is valid; \lstinline[style=tsinline]|Pick<User, 'missing'>| errors. This pattern is how utility types stay type-safe.
\end{solution}

\question[4]
Combining \lstinline[style=tsinline]|typeof| and \lstinline[style=tsinline]|keyof| as \lstinline[style=tsinline]|keyof typeof config| when \lstinline[style=tsinline]|config| is a \lstinline[style=tsinline]|const| object:
\begin{choices}
  \choice Always yields \lstinline[style=tsinline]|string|
  \CorrectChoice Yields a union of the object's keys (often narrow literal keys when \lstinline[style=tsinline]|as const| was used)
  \choice Deletes the object at compile time
  \choice Returns property values instead of keys
\end{choices}
\begin{solution}
\begin{lstlisting}[style=tsexam]
const config = { mode: 'dark', level: 1 } as const;
type ConfigKey = keyof typeof config;  // 'mode' | 'level'
\end{lstlisting}
\end{solution}

\question[4]
A \textbf{mapped type} \lstinline[style=tsinline]|{ [K in keyof User]: boolean }| describes:
\begin{choices}
  \choice An array of users
  \CorrectChoice An object type with the same keys as \lstinline[style=tsinline]|User| but every property type is \lstinline[style=tsinline]|boolean|
  \choice A function returning \lstinline[style=tsinline]|void|
  \choice A CSS module
\end{choices}
\begin{solution}
  Mapped types loop keys from \lstinline[style=tsinline]|keyof T|. Utilities like \lstinline[style=tsinline]|Partial| and \lstinline[style=tsinline]|Readonly| are implemented with this pattern in the standard library.
\end{solution}

\question[4]
Calling \lstinline[style=tsinline]|getProperty(user, 'email')| when \lstinline[style=tsinline]|email| is not a key of \lstinline[style=tsinline]|User|:
\begin{choices}
  \choice Succeeds because all strings are valid keys
  \CorrectChoice Fails type checking---the second argument must be \lstinline[style=tsinline]|keyof User|
  \choice Returns \lstinline[style=tsinline]|never| at runtime always
  \choice Requires \lstinline[style=tsinline]|as any| on \lstinline[style=tsinline]|user| only
\end{choices}
\begin{solution}
  \lstinline[style=tsinline]|K extends keyof T| ties the key parameter to the object type. This prevents typos like \lstinline[style=tsinline]|'emial'| when reading properties.
\end{solution}

\newpage

\uplevel{\par\textbf{Part 29: Readonly interfaces and immutability}\par}

\question[4]
In an interface, a \textbf{readonly} property:
\begin{choices}
  \choice May be reassigned freely after the object is created
  \CorrectChoice Cannot be assigned a new value after initialization (compile-time check on that property)
  \choice Forces the property to be optional
  \choice Exists only at runtime inside Angular templates
\end{choices}
\begin{solution}
\begin{lstlisting}[style=tsexam]
interface Point {
  readonly x: number;
  readonly y: number;
}
const p: Point = { x: 1, y: 2 };
// p.x = 5;  // compile error
\end{lstlisting}
  Readonly applies per property. Other non-readonly properties on the same object may still be mutable.
\end{solution}

\question[4]
The utility type \lstinline[style=tsinline]|Readonly<User>| compared with interface \lstinline[style=tsinline]|interface User { readonly id: number; name: string }|:
\begin{choices}
  \choice Are unrelated features
  \CorrectChoice \lstinline[style=tsinline]|Readonly<User>| makes \textbf{every} property of \lstinline[style=tsinline]|User| readonly in one step
  \choice \lstinline[style=tsinline]|Readonly<User>| deletes all optional fields
  \choice Only works on arrays, not interfaces
\end{choices}
\begin{solution}
  Hand-annotate important fields on an interface, or wrap an existing type with \lstinline[style=tsinline]|Readonly<T>| when passing config objects that must not be mutated (see Part~26).
\end{solution}

\question[4]
Assigning to \lstinline[style=tsinline]|obj.id| when \lstinline[style=tsinline]|obj| is typed \lstinline[style=tsinline]|Readonly<User>|:
\begin{choices}
  \choice Always succeeds because \lstinline[style=tsinline]|Readonly| is runtime-only
  \CorrectChoice Fails type checking if \lstinline[style=tsinline]|id| is readonly on that type
  \choice Requires \lstinline[style=tsinline]|as const| on every string field
  \choice Widens \lstinline[style=tsinline]|obj| to \lstinline[style=tsinline]|any|
\end{choices}
\begin{solution}
  \lstinline[style=tsinline]|Readonly| is shallow: nested objects may still be mutable unless their fields are also readonly or you use a deep readonly pattern.
\end{solution}

\question[4]
An interface with \lstinline[style=tsinline]|readonly items: string[]|:
\begin{choices}
  \choice Prevents \lstinline[style=tsinline]|push| on the array automatically (even though the element type is \lstinline[style=tsinline]|string[]|)
  \CorrectChoice Prevents reassigning \lstinline[style=tsinline]|items| to a new array; use \lstinline[style=tsinline]|readonly string[]| to block mutating methods like \lstinline[style=tsinline]|push|
  \choice Makes each string inside the array immutable at the character level automatically
  \choice Converts the interface to an \lstinline[style=tsinline]|enum|
\end{choices}
\begin{solution}
  Distinguish \textbf{readonly property} (cannot rebind \lstinline[style=tsinline]|items|) from \textbf{readonly array type} (cannot mutate array contents). Combine both when exposing internal lists safely.
\end{solution}

\question[4]
Why use a \textbf{readonly interface} for configuration passed into a function?
\begin{choices}
  \choice To speed up \texttt{tsc} by skipping checks
  \CorrectChoice To document that the callee must not mutate the caller's object (enforced at compile time)
  \choice To enable decorators automatically
  \choice To replace \lstinline[style=tsinline]|keyof|
\end{choices}
\begin{solution}
  Functions accepting \lstinline[style=tsinline]|Readonly<Config>| signal intent. Callers keep mutable objects; the parameter view is read-only. Runtime enforcement still needs discipline or \lstinline[style=tsinline]|Object.freeze| if required.
\end{solution}

\question[4]
Deep immutability for nested objects:
\begin{choices}
  \choice Is automatic for every \lstinline[style=tsinline]|Readonly<T>| without extra work
  \CorrectChoice Requires nested \lstinline[style=tsinline]|readonly| fields, recursive readonly types, or specialized utility types---\lstinline[style=tsinline]|Readonly<T>| is shallow only
  \choice Is provided by \lstinline[style=tsinline]|as const| on the outer object only in all cases
  \choice Cannot be expressed in TypeScript
\end{choices}
\begin{solution}
  Plan readonly boundaries at API edges. For deeply nested config, model nested interfaces with readonly fields or accept immutable data structures from the start.
\end{solution}

\uplevel{\par\textbf{Part 30: \texttt{as const} assertions}\par}

\question[4]
The assertion \lstinline[style=tsinline]|as const| on an object literal:
\begin{choices}
  \choice Widens every property type to \lstinline[style=tsinline]|string|
  \CorrectChoice Narrows properties to readonly literal types and marks the object readonly at the top level
  \choice Deletes optional properties
  \choice Emits a new enum at runtime
\end{choices}
\begin{solution}
\begin{lstlisting}[style=tsexam]
const routes = {
  home: '/',
  admin: '/admin',
} as const;
// routes.home is type '/' not string
\end{lstlisting}
  See also Part~21 for assertion basics.
\end{solution}

\question[4]
\lstinline[style=tsinline]|const colors = ['red', 'green', 'blue'] as const;| types \lstinline[style=tsinline]|colors| as:
\begin{choices}
  \choice \lstinline[style=tsinline]|string[]| mutable array
  \CorrectChoice A readonly tuple \lstinline[style=tsinline]|readonly ['red', 'green', 'blue']|
  \choice \lstinline[style=tsinline]|never[]|
  \choice \lstinline[style=tsinline]|number[]|
\end{choices}
\begin{solution}
  Without \lstinline[style=tsinline]|as const|, the type is \lstinline[style=tsinline]|string[]|. With it, lengths and element literals are fixed---useful for lookup tables and discriminant lists.
\end{solution}

\question[4]
\lstinline[style=tsinline]|as const| differs from typing a variable \lstinline[style=tsinline]|const x: string = 'left'| because:
\begin{choices}
  \choice They are always identical
  \CorrectChoice \lstinline[style=tsinline]|as const| preserves literal type \lstinline[style=tsinline]|'left'|; an explicit \lstinline[style=tsinline]|string| annotation widens to general \lstinline[style=tsinline]|string|
  \choice \lstinline[style=tsinline]|as const| removes readonly
  \choice \lstinline[style=tsinline]|as const| works only on classes
\end{choices}
\begin{solution}
  Prefer \lstinline[style=tsinline]|as const| on config objects; avoid annotating with a wide type unless you intentionally want widening.
\end{solution}

\question[4]
After \lstinline[style=tsinline]|const cfg = { mode: 'dark' } as const|, assigning \lstinline[style=tsinline]|cfg.mode = 'light'|:
\begin{choices}
  \choice Is always allowed
  \CorrectChoice Is a compile error (readonly literal property)
  \choice Changes the type of \lstinline[style=tsinline]|cfg| to \lstinline[style=tsinline]|any|
  \choice Requires \lstinline[style=tsinline]|Partial<cfg>|
\end{choices}
\begin{solution}
  Immutability is at the type level. Runtime mutation might still be possible on plain objects unless frozen---do not rely on \lstinline[style=tsinline]|as const| alone for security boundaries.
\end{solution}

\question[4]
\lstinline[style=tsinline]|keyof typeof routes| after \lstinline[style=tsinline]|const routes = { ... } as const| yields:
\begin{choices}
  \choice \lstinline[style=tsinline]|string| only
  \CorrectChoice A union of literal key names (for example \lstinline[style=tsinline]|'home' | 'admin'|)
  \choice The union of all route path string values only
  \choice \lstinline[style=tsinline]|never|
\end{choices}
\begin{solution}
  Pair \lstinline[style=tsinline]|as const| with \lstinline[style=tsinline]|typeof| and \lstinline[style=tsinline]|keyof| (Part~28) for type-safe keys without maintaining duplicate string union types.
\end{solution}

\question[4]
\lstinline[style=tsinline]|as const| is most appropriate when:
\begin{choices}
  \choice You need to mutate a lookup table at runtime frequently
  \CorrectChoice You want fixed literal values for keys, routes, action names, or discriminant tags in unions
  \choice You are casting unknown API data without validation
  \choice You are replacing \lstinline[style=tsinline]|interface| declarations with SQL
\end{choices}
\begin{solution}
  Action creators and Redux-style constants often use \lstinline[style=tsinline]|as const| so downstream \lstinline[style=tsinline]|switch| statements stay exhaustive.
\end{solution}

\uplevel{\par\textbf{Part 31: Type guards (narrowing in depth)}\par}

\question[4]
A \textbf{type guard} is:
\begin{choices}
  \choice A CSS selector for Angular components
  \CorrectChoice Code or a typed predicate that narrows a union to a more specific type in a branch (Part~8 introduces the idea)
  \choice A Karma test matcher only
  \choice The same as \lstinline[style=tsinline]|as any|
\end{choices}
\begin{solution}
  After a successful guard, the compiler knows which members are safe to access without assertion.
\end{solution}

\question[4]
\lstinline[style=tsinline]|if (typeof value === 'string')| inside a function where \lstinline[style=tsinline]|value: string | number|:
\begin{choices}
  \choice Leaves \lstinline[style=tsinline]|value| as the full union in the \texttt{if} block
  \CorrectChoice Narrows \lstinline[style=tsinline]|value| to \lstinline[style=tsinline]|string| in the \texttt{if} block and excludes \lstinline[style=tsinline]|string| in the \texttt{else} branch
  \choice Requires \lstinline[style=tsinline]|as string| on every use
  \choice Works only for objects, not primitives
\end{choices}
\begin{solution}
  \lstinline[style=tsinline]|typeof| guards suit primitives (\lstinline[style=tsinline]|string|, \lstinline[style=tsinline]|number|, \lstinline[style=tsinline]|boolean|, \lstinline[style=tsinline]|bigint|, \lstinline[style=tsinline]|symbol|, \lstinline[style=tsinline]|undefined|, \lstinline[style=tsinline]|function|). Remember \lstinline[style=tsinline]|typeof null === 'object'| at runtime.
\end{solution}

\question[4]
\lstinline[style=tsinline]|if (pet instanceof Dog)| narrows when:
\begin{choices}
  \choice \lstinline[style=tsinline]|pet| is an interface type only
  \CorrectChoice \lstinline[style=tsinline]|Dog| is a class constructor and \lstinline[style=tsinline]|pet| may be an instance of that class hierarchy
  \choice \lstinline[style=tsinline]|pet| is always \lstinline[style=tsinline]|null|
  \choice \lstinline[style=tsinline]|Dog| is a type alias for \lstinline[style=tsinline]|string|
\end{choices}
\begin{solution}
  \lstinline[style=tsinline]|instanceof| checks the prototype chain at runtime. It does not work for plain interface-shaped objects---use \lstinline[style=tsinline]|in| or discriminant fields instead.
\end{solution}

\question[4]
The \lstinline[style=tsinline]|in| operator guard \lstinline[style=tsinline]|if ('swim' in pet)| is useful when:
\begin{choices}
  \choice \lstinline[style=tsinline]|pet| is a number or string primitive
  \CorrectChoice Distinguishing plain object shapes in a union (for example \lstinline[style=tsinline]|Fish| vs \lstinline[style=tsinline]|Bird| interfaces)
  \choice Replacing \lstinline[style=tsinline]|strict| mode
  \choice Compiling templates to HTML
\end{choices}
\begin{solution}
  Combine with discriminated unions (\lstinline[style=tsinline]|kind| field) for the clearest narrowing. Part~18 covers discriminated unions; guards make them safe in \lstinline[style=tsinline]|if|/\lstinline[style=tsinline]|switch| bodies.
\end{solution}

\question[4]
A user-defined type guard \lstinline[style=tsinline]|function isUser(x: unknown): x is User|:
\begin{choices}
  \choice Must return \lstinline[style=tsinline]|User| at runtime always
  \CorrectChoice Returns a boolean; when \texttt{true}, TypeScript treats \lstinline[style=tsinline]|x| as \lstinline[style=tsinline]|User| in that branch
  \choice Replaces runtime validation entirely
  \choice Only works inside enums
\end{choices}
\begin{solution}
  You are responsible for correct runtime checks (property presence, types). Wrong predicates can lie to the compiler and cause runtime errors later.
\end{solution}

\question[4]
An \textbf{assertion function} \lstinline[style=tsinline]|function assertIsNumber(x: unknown): asserts x is number|:
\begin{choices}
  \choice Returns \lstinline[style=tsinline]|boolean| without throwing
  \CorrectChoice Throws or narrows---if the function returns normally, \lstinline[style=tsinline]|x| is treated as \lstinline[style=tsinline]|number| afterward
  \choice Is the same as \lstinline[style=tsinline]|as number| always
  \choice Only applies to class fields
\end{choices}
\begin{solution}
  Useful in test helpers and validation layers. Unlike a predicate, failed assertions do not enter the narrowed branch---they exit via throw.
\end{solution}

\question[4]
Checking \lstinline[style=tsinline]|if (value === null)| when \lstinline[style=tsinline]|value: string | null|:
\begin{choices}
  \choice Leaves \lstinline[style=tsinline]|value| as \lstinline[style=tsinline]|string | null| in the true branch
  \CorrectChoice Narrows out \lstinline[style=tsinline]|null| in the false branch (and to \lstinline[style=tsinline]|null| in the true branch)
  \choice Requires \lstinline[style=tsinline]|as const|
  \choice Works only with \lstinline[style=tsinline]|Readonly<T>|
\end{choices}
\begin{solution}
  Equality checks (\lstinline[style=tsinline]|===|, \lstinline[style=tsinline]|!==|) and truthiness checks narrow unions. With \texttt{strictNullChecks}, handle both \lstinline[style=tsinline]|null| and \lstinline[style=tsinline]|undefined| explicitly when needed.
\end{solution}

\endgradingrange{csawebts}
